% BFZI-first reconstruction (2026-10-31). XeLaTeX + BibTeX.
\documentclass[UTF8,11pt,a4paper]{ctexart}
\usepackage{amsmath,amssymb,mathtools,bm}
\usepackage{booktabs,longtable,geometry,hyperref}
\geometry{margin=25mm}
\hypersetup{hidelinks}
\newcommand{\jj}{\mathrm j}
\newcommand{\RR}{\mathbb R}
\newcommand{\Herm}[1]{\operatorname{Herm}\left(#1\right)}
\newcommand{\grad}{\nabla}
\newcommand{\Div}{\nabla\!\cdot}
\newcommand{\Curl}{\nabla\!\times}
\title{BFZI 连续螺旋水下无线电能传输\\电磁--热--谐振电路联合代理\\\large 新理论架构（研究草案）}
\author{SDFMPNEO 理论重构}
\date{2026年11月6日（物理直接法 V1.0 正式锁定；G1 仍未验收）}
\begin{document}
\maketitle
\begin{abstract}
本稿取代 SDFMPNEO 旧理论，以最新版 0UnifiedModelingUWPT 所建立的圆形--圆角方形连续螺旋
及跨匝仿射结构为几何基础，研究含独立封装、海水导电反作用、六自由度位姿、空间电磁损耗、
谐振电路与连续时间热动力学的统一代理架构。解析场积分、受约束非局部神经响应、
空间损耗一致性和守恒热降阶为候选技术。尚未通过实验验证的假设不视为既有能力。
\end{abstract}
\tableofcontents
\section{范围、来源与物理约定}
\textbf{几何权威来源：}\href{https://github.com/AliceThunder/0UnifiedModelingUWPT}{0UnifiedModelingUWPT} 的
\texttt{sections/02\_geometry\_model.tex}。
海水源项与导电反作用参考 \texttt{03\_package\_seawater\_impedance.tex}，
匝号结构化快速计算参考 \texttt{04\_geometry\_induced\_fast\_computation.tex}。
本文不沿用旧的双阶数超椭圆模型；“统一建模”严格指圆形与圆角方形之间的
\textbf{连续平面螺旋几何统一}，封装空腔只是海水导电域的结构修正。

研究目标是从同一场景及其运行历程在合法查询点生成
\(\bm E,\bm B,\bm J,q_i,T_i\)，在端口生成复阻抗矩阵 \(\bm Z\)，
在热时间上实现 \(T\to\mathrm{material}\to\bm Z\to\bm I\to q\to T\) 的双向反馈。
仅预测测温点、仅预测 \(R,L,M\) 或仅单向将损耗传给热场均不足以达成完整目标。

使用峰值相量 \(e^{\jj\omega t}\)，时间平均功率前系数为 \(1/2\)；
\(\Herm{\bm Z}=(\bm Z+\bm Z^{\mathrm H})/2\)。
首版采用磁准静态、线性各向同性海水、非导电封装和明确补偿拓扑，
非线性磁芯、各向异性介质、快速运动与全波问题须分阶段扩展验证。
无固定世界网格推理不等于无需局部求积、数值参考或 FEM 教师数据。
现有 SDFMPNEO 的 \texttt{SuperellipseSpiral} 也不等于 BFZI 最新几何类，不得直接重命名冒充迁移。
\section{最新版 BFZI 圆形--圆角方形连续螺旋}\label{sec:geometry}
\subsection{统一绕行参数和尺度}
依照 0UnifiedModelingUWPT 的原始公式，定义
\begin{equation}
\eta=k+\frac{\tau}{2\pi},\quad
\beta=\frac{s}{2\pi},\quad
\psi=\sigma_w\tau,\quad
k=0,\ldots,N-1,\ \tau\in[0,2\pi].
\end{equation}
其中 \(\sigma_w=\pm1\) 指定绕制方向。
直边外侧半尺寸、角部尺度及直边半长度为
\begin{equation}
\boxed{\begin{aligned}
A(\eta)&=A_0-s\eta,\\
\varrho(\eta)&=\varrho_0-\alpha s\eta,\\
\chi(\eta)&=A_0-\varrho_0-(1-\alpha)s\eta,
\end{aligned}}\qquad
A=\chi+\varrho,\quad 0\le\alpha\le1.
\label{eq:bfzi-scales}
\end{equation}
\(s\) 是直边处的标称匝距；仅当 \(\alpha=1\) 时，圆角部分也保持恒定法向匝距。
合法参数要求 \(A>0,\varrho>0,\chi\ge0\)、线元不退化，
且实际导体/绝缘截面满足净距、自交、封装约束。

\subsection{四组直边和四组角部}
定义 \(\bm n(\psi)=(\cos\psi,\sin\psi)^{\mathsf T}\) 和
\(\bm t(\psi)=(-\sin\psi,\cos\psi)^{\mathsf T}\)。
第 \(j=0,\dots,3\) 个直边在
\(\eta_{j,k}=k+j/4,\psi_j=\sigma_wj\pi/2,\xi\in[-1,1]\) 上满足
\begin{equation}
\begin{aligned}
\bm r_{j,k}(\xi)&=A(\eta_{j,k})\bm n(\psi_j)
 +\sigma_w[\chi(\eta_{j,k})\xi-\beta]\bm t(\psi_j),\\
\bm g_{j,k}(\xi)&=\sigma_w\chi(\eta_{j,k})\bm t(\psi_j).
\end{aligned}
\label{eq:bfzi-straight}
\end{equation}
第 \(m=0,\dots,3\) 个角部在
\(\tau_m=m\pi/2+u,\eta_{m,k}=k+\tau_m/(2\pi)\)，
\(\psi_m=\sigma_w\tau_m,\ u\in[0,\pi/2]\)，
\(\bm e_m=\bm n(\sigma_wm\pi/2)+\sigma_w\bm t(\sigma_wm\pi/2)\) 下，
通过支撑函数
\[
h_{m,k}(u)=\chi(\eta_{m,k})\bm e_m\cdot\bm n(\psi_m)
+\varrho(\eta_{m,k})
\]
产生实际中心线和有向线元
\begin{equation}
\begin{aligned}
\bm r_{m,k}(u)&=h_{m,k}\bm n(\psi_m)
 +\sigma_wh'_{m,k}(u)\bm t(\psi_m),\\
\bm g_{m,k}(u)&=\sigma_w[
 \varrho(\eta_{m,k})
 -2\sigma_w(1-\alpha)\beta\,\bm e_m\cdot\bm t(\psi_m)]
 \bm t(\psi_m).
\end{aligned}
\label{eq:bfzi-corner}
\end{equation}
两式的 \(\bm g=\partial_{\zeta}\bm r\) 都是有向线元向量；
角部与前后直边、各匝之间连续连接成一根从外端到内端的导体。
\textbf{不应将其替换为各匝互不连接的圆环}。

\subsection{圆形退化及跨匝仿射关系}
当 \(\varrho_0=A_0\) 且 \(\alpha=1\)，有 \(\chi=0\)；
四个角部共同构成圆形连续平面螺旋
\[
\bm r_{\rm cir,k}(\tau)
=A(\eta)\bm n(\sigma_w\tau)-\sigma_w\beta\bm t(\sigma_w\tau).
\]
切向偏移项不能删去。记直边/角部线段组 \(b=0,\dots,7\)，
固定各组局部参数 \(\zeta_b\) 时，原模型严格满足
\begin{equation}
\boxed{
\bm r_{b,k}=\bm r_{b,0}+k\Delta\bm r_b,\qquad
\bm g_{b,k}=\bm g_{b,0}+k\Delta\bm g_b.
}
\label{eq:bfzi-affine}
\end{equation}
这是后续快速积分和代理场源编码的几何基础，并非神经网络训练得来的近似。

\subsection{独立六自由度装配}
令两侧组件 \(p\) 具有 \(R_p\in SO(3)\) 和
\(\bm c_p\in\mathbb R^3\)，则
\begin{equation}
\bm r_{p,b,k}^{(3)}
 =\bm c_p+R_p[\bm r_{b,k};0],\qquad
\bm g_{p,b,k}^{(3)}=R_p[\bm g_{b,k};0].
\end{equation}
刚体变换保持跨匝仿射结构。
相对姿态可表示为 \(R_1^{\mathsf T}R_2\)，相对位置为
\(R_1^{\mathsf T}(\bm c_2-\bm c_1)\)。

\section{封装结构、海水空间与任意位姿场景}\label{sec:scene}
\subsection{真实材料域}
每个刚性装配件含 BFZI 线圈、导体截面、绝缘、灌封材料、
磁芯、屏蔽和独立参数化的封装盒外壳（可按工况选配）。
盒体的尺寸与壁厚不由 \(A_0\) 唯一决定，
盒体和线圈共享或分别指定合法的位姿。
实体导热区域记为 \(\Omega_i(t)\)，不允许不同材料区域非法重叠；
对气密封装，其海水域可定义为
\[
\Omega_s(t)=\mathbb R^3\setminus
\bigcup_p\overline{\Omega_{\mathrm{excluded},p}(t)}.
\]
这不是无限空间离散网格，而是几何定义；数值参考可以采取
人工截断域、无穷远边界或边界积分并进行收敛研究。

\subsection{几何材料查询与接口}
给每个组件局部坐标 \(\bm\xi_p=R_p^{\mathsf T}(\bm x-\bm c_p)\)，
材料边界可由解析几何、参数曲面或符号距离函数
\(d_i(\bm\xi_p;\bm g_i)\) 查询。
\textbf{SDF 是分区和求积辅助，不取代 BFZI 原始螺旋中心线。}
在非导电封装/海水界面，海水法向导电电流为零；
导电盒壁与海水相邻时必须使用真实电磁界面连续性，
不可套用绝缘条件。热学还必须区分固体内部导热、
接触传热及盒壁--海水热交换。

\subsection{材料范围和海水运动}
电磁材料可依赖 \(\bm x,T,\omega\)，包括
\(\sigma,\mu,\epsilon\)；热材料包括 \(\rho,c,k\)，
可分段均匀或在参考求解器允许时变参。
首版建议均匀各向同性海水、线性材料、
非导电非磁性封装；导电/磁性盒体、频散、非线性、
各向异性和任意空间变化是分阶段扩展，不能未经验证声称全覆盖。
世界坐标的海流 \(\bm v_s(\bm x,t)\)、重力和外界背景
不随组件的局部坐标旋转。任意\textbf{静态姿态}与高速运动不同：
若运动电动势、流固边界变化显著，应增加移动介质的电磁/热方程。

\section{BFZI 解析电磁参考与海水非局部反作用}\label{sec:em}
\subsection{完整端口回路及空气电感}
对单位端口电流 \(p\)，将主螺旋及明确的馈电/回流段统一记为
\(\Gamma_p\)，定义
\[
\bm C_p(\bm x)=\frac1{4\pi}
\int_{\Gamma_p}\frac{\mathrm d\bm\ell'}{\|\bm x-\bm r'\|}.
\]
对主螺旋在固定 Gauss 节点 \(b,\mu\) 的带权线元，
\(\bm r_k=\bm r_0+k\Delta\bm r\)，
\(\bm h_k=\bm h_0+k\Delta\bm h\)，
可对 \(\sum_k\bm h_k/\|\bm x-\bm r_k\|\)
按 BFZI 原稿采用特殊情形精确求和、
一般情形误差受控二阶分块，近奇异时回退直接计算。
非重复馈/回流段单独积分，不可遗漏。

空气电感参考以 Neumann 积分为基础
\[
L^{\rm air}_{pq}=\frac{\mu_0}{4\pi}
\int_{\Gamma_p}\int_{\Gamma_q}
\frac{\mathrm d\bm\ell_p\cdot\mathrm d\bm\ell_q}
{\|\bm r_p-\bm r_q\|}.
\]
同端口自项必须考虑有限导体截面/GMD 和近奇异局部积分，
不能直接对零半径自重合线元求积。
双匝号节点对令
\(Q_{kl}=\|\bm d_0+k\Delta\bm x-l\Delta\bm y\|^2\) 和
\(H_{kl}=(\bm h_0+k\Delta\bm h)\cdot(\bm z_0+l\Delta\bm z)\)；
\(Q\) 为二次多项式，\(H\) 分别关于 \(k,l\) 至多一次。
当两匝号增量相等或相反时可精确转为匝号差/和的单索引求和，
否则以块矩和余项界做加速。
该结构已经来自几何论文，不应说是新神经算子自己发明的。

\subsection{海水导电反作用}
在峰值相量 \(e^{\jj\omega t}\) 下使用磁准静态 \(A\)--\(\phi\)：
\begin{equation}
\bm E=-\jj\omega\bm A-\grad\phi,\quad
\bm B=\Curl\bm A,\quad
\bm J_s=\sigma_s\bm E,\quad \Div\bm J_s=0.
\end{equation}
海水与绝缘封装面满足 \(\bm n\cdot\bm J_s=0\)。
在源仓库相应的分段材料与离散前提下，场的自洽反作用可写作
\begin{equation}
\boxed{(\mathsf I+\mathcal KP_\phi)\bm X_p=\bm C_p},
\label{eq:sea-response}
\end{equation}
其中 \(P_\phi\) 实现电流连续约束，\(\mathcal K\) 描述涡流
反过来对磁场的影响。统一得到端口海水反作用复阻抗
\(\bm Z_{\rm sw}\)，并与 \(j\omega\bm L^{\rm air}\)
组成源文献的复阻抗基底。这个基底的实部指海水导电反作用，
\textbf{不含导体自身欧姆、趋肤、邻近损耗}。

\subsection{扩展求解器和高精度参考}
非导电封装基底之外，导电外壳、磁芯和复杂介质需要适用的
材料边界/弱式、界面积分或体积分，不得假定式
\eqref{eq:sea-response} 适用于任意非线性材料。
为每个端口保存同源复电磁状态，允许同时读出
\(\bm E,\bm J,\bm B\)、端口电压及空间材料损耗，
作为后续神经学习的物理教师。

\section{新推导：保结构海水涡流 Galerkin 与同源端口阻抗}\label{sec:galerkin}
\subsection{问题范围及导电电流函数空间}
本节\textbf{首次将已有 BFZI 解析场源与新的保结构海水电流低阶模型相结合}。
限定如下：谐波峰值相量 \(e^{\jj\omega t}\)、磁准静态、线性互易介质、
均匀 \(\mu=\mu_0\)、海水电导率 \(0<\sigma_{\min}\leq\sigma_s(\bm x)\leq\sigma_{\max}<\infty\)、
绝缘封装及海水外部的正确衰减条件，端口回路与海水界面不相交。
所有基函数和积分取实值几何/材料，复数仅由谐波场系数引入。
静态位姿可任意合法改变，但暂不纳入运动感应电动势。
对可计算有限域 \(\Omega_s\)（或合适的有限能量无限域空间），定义
\begin{equation}
\mathcal V(\Omega_s)
=\{\bm v\in H(\mathrm{div};\Omega_s):
\nabla\cdot\bm v=0,\ 
\bm n\cdot\bm v|_{\partial\Omega_s}=0\quad
\text{（有限域中含人工绝缘外边界）}\}.
\label{eq:sea-solenoidal-space}
\end{equation}
对无限海水，还需要无穷远处衰减、能量有限及适当的函数空间闭合；
后文所有有限维证明先在边界良定、积分收敛的条件下成立。
\textbf{多连通域中仅使用局部 curl 基可能遗漏拓扑环流自由度}；
基空间必须含必要的全局回路模态或采用等价的连续性投影。
这是一项待落实的具体离散构造，不可把其存在性等同于已完成快速实现。

令 \(\bm W_m(\bm x;\mathcal G)\in\mathcal V\) 为 \(r\) 个线性独立的
实向量基，封装与位姿以 \(\mathcal G\) 表示；对任意端口电流列
\(\bm I\in\mathbb C^P\)，采用
\begin{equation}
\bm J_s(\bm x)=\sum_{m=1}^r z_m\bm W_m(\bm x)
=:\bm W(\bm x)\bm z .
\label{eq:sea-reduced-current}
\end{equation}
\(\bm W\) 的列代表电流密度模式，\(\bm z\) 为复数电流模式系数。
允许 \(\bm W\) 随几何变化，但在每一次线性频域求解中应作为固定空间；
随温度改变的 \(\sigma_s\) 进入下面的 \(\bm R_s\)。

\subsection{从 \(A\)--\(\phi\) 方程得到约束弱式}
定义自由空间积分核
\begin{equation}
(\mathcal G_0\bm f)(\bm x)
=\int_{\Omega_s}\frac{\bm f(\bm y)}
{4\pi\|\bm x-\bm y\|}\,\mathrm d\bm y,
\label{eq:free-g}
\end{equation}
以及由 BFZI 完整端口路径（含明确馈/回流）形成的单位端口线源
\(\bm C_p(\bm x)=\int_{\Gamma_p}(4\pi\|\bm x-\bm r'\|)^{-1}
\,\mathrm d\bm\ell'\)。
在 \(\mu=\mu_0\) 条件下，
\[
\bm A(\bm x)
=\mu_0\sum_{p=1}^P \bm C_p(\bm x)I_p
+\mu_0\mathcal G_0\bm J_s(\bm x),\quad
\sigma_s^{-1}\bm J_s=-\jj\omega\bm A-\nabla\phi .
\]
取任意 \(\bm v\in\mathcal V\) 测试，上式与 \(\bm v\) 做
\(L^2\) 内积并对 \(\Omega_s\) 积分。
对采用人工绝缘外边界的\emph{有限域}，
由分部积分、\(\nabla\cdot\bm v=0\) 和\emph{整个边界}
的测试场零法向迹，可得
\(\int_{\Omega_s}\bm v\cdot\nabla\phi\,\mathrm dV=0\)。
若人工外边界允许非零法向通量，则该边界项
不可直接消去；详见第 \ref{sec:g1-foundation-recheck} 节。
故有
\begin{equation}
\int_{\Omega_s}\sigma_s^{-1}\bm v\cdot\bm J_s\,\mathrm dV
+\jj\omega\mu_0\int_{\Omega_s}\bm v\cdot\mathcal G_0\bm J_s\,\mathrm dV
=-\jj\omega\mu_0\sum_p I_p\int_{\Omega_s}\bm v\cdot\bm C_p\,\mathrm dV .
\label{eq:sea-weak}
\end{equation}
采用 \(\bm v=\bm W_m\) 形成有限维矩阵
\begin{align}
[\bm R_s]_{mn}
&=\int_{\Omega_s}\sigma_s^{-1}\bm W_m\cdot\bm W_n\,\mathrm dV,
\label{eq:Rs}\\
[\bm L_s]_{mn}
&=\mu_0\int_{\Omega_s}\bm W_m\cdot\mathcal G_0\bm W_n\,\mathrm dV,
\label{eq:Ls}\\
[\bm K]_{mp}
&=\mu_0\int_{\Omega_s}\bm W_m\cdot\bm C_p\,\mathrm dV .
\label{eq:Kbfzi}
\end{align}
其中 \(\bm K\) 并非数学刚度矩阵；此处专用作
\textbf{海水电流模式--BFZI 端口互感耦合矩阵}。
上述基于实值模式的 \(\bm R_s,\bm L_s,\bm K\) 均为实矩阵，
其中 \(\bm R_s\) 对称正定，\(\bm L_s\) 在适用自由空间核和电流空间中对称半正定，
其物理含义分别为电阻功率与磁场能量度量。
令
\(\bm A_r(\omega)=\bm R_s+\jj\omega\bm L_s\)，
得到
\begin{equation}
\boxed{\bm A_r\bm z=-\jj\omega\bm K\bm I,\qquad
\bm z=-\jj\omega\bm A_r^{-1}\bm K\bm I.}
\label{eq:sea-current-state}
\end{equation}
因为任意非零 \(\bm v\in\mathbb C^r\) 有
\(\operatorname{Re}(\bm v^{\mathrm H}\bm A_r\bm v)
=\bm v^{\mathrm H}\bm R_s\bm v>0\)，
\(\bm A_r\) 在实频率上非奇异。实际采用适当频率与材料归一化，
避免 \(\bm R_s\) 和 \(\omega\bm L_s\) 尺度相差巨大造成数值条件恶化。

\subsection{端口散射阻抗与同源功率定理}
由自由空间 Green 核的交换对称性和相同端口/回路方向，
涡流磁矢势对第 \(p\) 个端口产生的附加链磁通为
\(\lambda_{s,p}=\sum_m[\bm K]_{mp} z_m\)。
按 \(+\jj\omega\lambda\) 的端口感应电压约定，
\(\bm v_s=\jj\omega\bm K^{\mathsf T}\bm z\)。
代入式 \eqref{eq:sea-current-state} 得到
\begin{equation}
\boxed{
\bm Z_{\rm sw}(\omega)
=\omega^2\bm K^{\mathsf T}
(\bm R_s+\jj\omega\bm L_s)^{-1}\bm K .}
\label{eq:zsw-schur}
\end{equation}
\textbf{定理（有限维互易、被动与海水功率一致性）}：
在本节所有假设成立、\(\bm R_s=\bm R_s^{\mathsf T}\succ0\)、
\(\bm L_s=\bm L_s^{\mathsf T}\succeq0\)、\(\bm K\) 实值，
且响应、端口和局部损耗均采用\textbf{同一} \(\bm W\)
和矩阵 \eqref{eq:Rs}--\eqref{eq:Kbfzi} 的条件下，
有
\begin{equation}
\bm Z_{\rm sw}=\bm Z_{\rm sw}^{\mathsf T},\quad
\operatorname{Herm}(\bm Z_{\rm sw})\succeq0,
\label{eq:sea-passivity-theorem}
\end{equation}
并对任意复端口电流 \(\bm I\) 满足
\begin{equation}
\boxed{
\frac12\bm I^{\mathrm H}\operatorname{Herm}(\bm Z_{\rm sw})\bm I
=\frac12\bm z^{\mathrm H}\bm R_s\bm z
=\int_{\Omega_s}q_s(\bm x)\,\mathrm dV,
\quad
q_s=\frac{|\bm W(\bm x)\bm z|^2}{2\sigma_s(\bm x)}.
}
\label{eq:sea-power-proof}
\end{equation}
证明：\(\bm A_r^{\mathsf T}=\bm A_r\) 导致
\((\bm A_r^{-1})^{\mathsf T}=\bm A_r^{-1}\)，故阻抗对称。
由
\[
\operatorname{Herm}(\bm A_r^{-1})
=\bm A_r^{-\mathrm H}\bm R_s\bm A_r^{-1}
\]
可得
\[
\operatorname{Herm}(\bm Z_{\rm sw})
=\omega^2\bm K^{\mathsf T}
\bm A_r^{-\mathrm H}\bm R_s\bm A_r^{-1}\bm K\succeq0.
\]
结合 \(\bm z=-\jj\omega\bm A_r^{-1}\bm K\bm I\)
得到前两个功率表达相等；将 \eqref{eq:Rs}
和 \eqref{eq:sea-reduced-current} 代入则得到体积分表达。
证毕。

该定理只证明\textbf{同一降阶模型内部}的物理一致性，
不证明 \(\bm W\) 足以逼近真实海水涡流；
积分数值误差若未一致处理也会破坏数值等式。
在 \(\omega\to0\) 且矩阵适当有界时，
\(\bm Z_{\rm sw}=\omega^2\bm K^{\mathsf T}\bm R_s^{-1}\bm K
+O(\omega^3)\)，提供低频渐近检验。

\subsection{热投影可直接取自同一电流系数}
对任意固体/海水热测试基 \(\varphi_a(\bm x)\)，
海水损耗投影为
\[
b_a^{\rm sea}
=\int_{\Omega_s}\varphi_a(\bm x)
\frac{|\bm W(\bm x)\bm z|^2}{2\sigma_s(\bm x)}\,\mathrm dV
=\frac12\bm z^{\mathrm H}\bm Q_a\bm z ,
\quad
[\bm Q_a]_{mn}=
\int_{\Omega_s}\frac{\varphi_a\bm W_m\cdot\bm W_n}{\sigma_s}\,\mathrm dV.
\]
此式建立同一涡流状态到\textbf{热空间输入}的明确接口。
若热测试基含常数分量 \(\varphi_0\equiv1\)，
其对应 \(\bm Q_0=\bm R_s\)，可直接满足总能量约束。
温度场和空间热源均可按合法场景在任意查询点继续解码，
但复杂海水外部流动应另由热/流动算子处理。

\section{几何变化下合规海水电流空间的具体构造}\label{sec:admissible-atlas}
\subsection{研究任务和必要的限制}
第 \ref{sec:galerkin} 节要求海水电流基
\(\bm W_m\in H(\mathrm{div};\Omega_s)\) 且
\(\nabla\cdot\bm W_m=0\)、
\(\bm n\cdot\bm W_m=0\)（绝缘实体界面），
第 \ref{sec:neural-rom} 节仅假设这类字典已存在。
现在将其变为\textbf{可实施的构造路线}。
限制在封闭、相互不交的非导电封装体，
表面至少分片 \(C^1\)，最小间隙 \(h_{\min}>0\)，
有限外域或能量收敛的无穷海水域。
完整的任意拓扑、尖角、近接触极限和移动导电接触将在单独图谱中处理。
\textbf{“任意静态六自由度位姿”仅指满足非碰撞、正间隙与局部形状图谱条件的工况。}

\subsection{构造 I：边界消失矢量势的 curl 基}
在海水中设 \(\mathcal U_\alpha\) 是局部求积/表示图谱，
\(\chi_\alpha\in C^1_c(\mathcal U_\alpha)\) 是有限支撑窗函数；
局部封装面由
\(d_\alpha(\bm x)=0\) 定义，
其中 \(d_\alpha>0\) 在海水侧，且
\(\nabla d_\alpha\ne0\) 于光滑界面。
用分片光滑向量函数 \(\bm a_{\alpha m}\) 构造
\begin{equation}
\boxed{
\bm V_{\alpha m}(\bm x)
=d_\alpha(\bm x)\chi_\alpha(\bm x)\bm a_{\alpha m}(\bm x),
\qquad
\bm W_{\alpha m}(\bm x)=\nabla\times\bm V_{\alpha m}(\bm x).
}
\label{eq:curl-atlas}
\end{equation}
将 \(\mathcal U_\alpha\) 的支撑限制在
\(\Omega_s\) 内及其当前指定的绝缘界面附近，
且不跨越其它固体。
于是恒等式 \(\nabla\cdot\nabla\times\bm V=0\)
在分布意义下保证连续性。
由于 \(\bm V_{\alpha m}=0\) 于相关绝缘界面，
其切向迹为零；由
\(\bm n\cdot(\nabla\times\bm V)=
\operatorname{div}_{\Gamma}(\bm V\times\bm n)\)
可得
\begin{equation}
\boxed{\nabla\cdot\bm W_{\alpha m}=0,\qquad
\bm n\cdot\bm W_{\alpha m}=0.}
\label{eq:curl-constraints}
\end{equation}
更一般地只需 \(\bm n\times\bm V=0\)，
不必强制矢量势的全部分量消失。
在界面 \(d_\alpha=0\) 上，
\(\nabla\times(d_\alpha\chi_\alpha\bm a)
=\chi_\alpha\nabla d_\alpha\times\bm a\)，
这是切向非零的涡流模式，\textbf{并不会把真实海水界面切向电流错误设为零}。
若支撑完全位于海水内部，则无需上述边界消失因子，
直接取 \(\bm W=\nabla\times(\chi\bm a)\)。

对于有尖角的盒壁，距离函数可能只具 Lipschitz 正则性，
需使用分片 \(H(\mathrm{curl})\) 势及合法迹，
或将光滑面片分区并用相容拼接；不能在不可微的
\(\min(d_1,d_2)\) 上盲目应用经典点态 curl 公式。
可采用解析参数面、圆角矩形盒体图谱等，
\textbf{不得将用于界面图谱的圆角盒体函数误替代 BFZI 线圈中心线定义}。

\subsection{构造 II：局部 Piola 图谱和刚体复用}
对于可用 \(C^1\) 微分同胚
\(F_{\alpha,\bm g}:\widehat{\mathcal U}_\alpha
\to\mathcal U_\alpha(\bm g)\) 描述的局部几何形变，
令 \(\bm J_F=D_{\widehat{\bm x}}F_{\alpha,\bm g}\)，
\(\det\bm J_F>0\)；定义\textbf{逆变} Piola 场
\begin{equation}
\boxed{
\bm W_{\alpha m}(\bm x;\bm g)
=\left.\frac{\bm J_F\,\widehat{\bm W}_{\alpha m}}
{\det\bm J_F}\right|_{\widehat{\bm x}=F^{-1}_{\alpha,\bm g}(\bm x)}.
}
\label{eq:piola-atlas}
\end{equation}
它满足
\[
\nabla_{\bm x}\cdot\bm W_{\alpha m}
=\frac{\widehat\nabla\cdot\widehat{\bm W}_{\alpha m}}
{\det\bm J_F},
\]
并保持界面的法向通量；
因此参考图谱中符合约束的电流模式可迁移到变化的封装几何。
\textbf{Piola 映射是既有数学工具，非本研究原创}；
其具体优势在于几何变化后不必重新学习散度约束。

纯刚体运动下 \(F(\widehat{\bm x})=R\widehat{\bm x}+\bm c\)，
\(\bm J_F=R\)、\(\det R=1\)，故
\(\bm W(\bm x)=R\widehat{\bm W}(R^\top(\bm x-\bm c))\)。
在均匀海水中，
该\textbf{同一局部模式}的电阻内积、自磁能和自身几何核保持刚体不变；
不同封装之间的\textbf{交叉磁能/间隙场}仍依赖相对位姿，
必须更新，不得沿用每盒独立求和。

\subsection{双封装图谱：近盒、间隙和远场}
只对单一封装建立参考域 Piola 图谱不足以覆盖
双封装六自由度互相错位后的海水区域，
且一般不存在一个具有统一低形变界的
\textbf{全参数域全局微分同胚}。
建议将字典分为：
\begin{equation}
\boxed{
\mathcal W(\mathcal G)
=\mathcal W_{\mathrm{near},1}
+\mathcal W_{\mathrm{near},2}
+\mathcal W_{\mathrm{gap}}
+\mathcal W_{\mathrm{far}}
+\mathcal W_{\mathrm{top}} .
}
\label{eq:dictionary-decomp}
\end{equation}
近盒图谱随各自装配件刚体运动或局部尺寸伸缩；
两盒靠近时增加随正间隙 \(h\) 缩放的间隙图谱；
远场采用合适衰减的解析模式、非局部边界映射或自适应局部模式；
多连通电流域需要相应的\textbf{全局拓扑谐和环流}，否则 curl 字典不完备。
这里的加号表示有限基空间之和，\textbf{不是}将每个盒子
\emph{独立求解的}阻抗相加。

采用近盒局部基时，需保证其支撑不侵犯第二个盒体；
相对姿态变化使其发生交叠时，动态裁剪或用新的
\(\chi_\alpha\) 重新生成合规模式，并重新计算受影响的 Gram 项。
间隙 \(h\to0^+\) 时，Piola Jacobian 条件数和局部场梯度可能恶化，
应增广近场模式、缩小适用范围或回退参考求解；
不宣称固定低阶数在所有近接触极限仍然准确。

\subsection{几何网络应学习什么}
网络仅学习以下\textbf{连续空间的选择/压缩}：
\[
\bm S_\theta:
(A_0,\varrho_0,s,\alpha,N,\sigma_w,\bm g_{\rm box},
q_{\rm rel},\omega,\sigma_s,\text{thermal summary})
\mapsto\mathbb R^{M\times r}.
\]
约束模式字典的散度/法向条件与网络权重无关，
低维子空间 \( \bm W_\theta=\bm W_0\bm S_\theta\)
仍合规。用局部图谱、特征缩放、秩条件以及
必要的局部正交化维持 \(\bm R\succ0\)；
随几何改变导致参考字典基数量变化时，
使用固定上界模式槽/掩码或可变大小集合编码器，
并明确模型接口。若新场景的模式空间覆盖不足，
应增加图谱或教师样本，而不是简单加大学习率。

\subsection{推理阶段无需世界 FEM 网格，但不能免积分}
几何通过解析图谱、局部积分节点和 BFZI 结构化线源处理，
不必构造固定世界 FEM 网格。体积分和核积分仍然存在；
每个图谱的积分和映射必须进入计算复杂度预算。
本节证明的是\textbf{约束保持}，
不是字典足够完备、跨几何高精度或在线快速的证明。

\section{BFZI 结构化场源与受约束神经响应}\label{sec:surrogate}
\subsection{从端口黑箱回归转为同源场代理}
仅独立拟合 \(R,L,M\) 或损耗、温度容易破坏
电流连续性、互易性、端口功率与空间热源的一致性。
因此研究同一个复场响应（或其低维状态）的学习方案，
再从同一响应读取端口和空间量。
用 \(\bm f_{\rm BFZI}\) 表示跨匝仿射几何节点、
聚合后的线源矩和已知空气积分先验；
用 \(\mathcal G\) 表示封装几何、材料分区、海水边界与相对位姿。
假设可建立低成本物理先验 \(\bm X_p^{(0)}\)，则定义候选
\begin{equation}
\boxed{
\widehat{\bm X}_p(\bm x)
=\bm X_p^{(0)}(\bm x)+
\mathcal P_{\mathcal G}
\mathcal N_\theta(\bm f_{{\rm BFZI},p},\mathcal G,\omega,\bm a;\bm x).
}
\end{equation}
其中 \(\bm a\) 为慢时间热状态，\(\mathcal P_{\mathcal G}\)
表示与材料、连续性和界面一致的投影/闭合结构。
该算子尚未构成完整证明；需要比较散度自由基、
低维势约束、边界积分和少量局部求积等具体实现，
并将 \(\mathcal P_{\mathcal G}\) 的计算成本计入推理时间。

\subsection{空间场、阻抗和损耗同源读取}
相量场由相应 \(A,\phi\) 读出
\[
\widehat{\bm B}=\Curl\widehat{\bm A},\quad
\widehat{\bm E}=-\jj\omega\widehat{\bm A}-\grad\widehat\phi,\quad
\widehat{\bm J}_i=\sigma_i\widehat{\bm E}_i.
\]
实际使用的 \(\bm X\) 定义和规范必须与源物理算子一致；
端口矩阵和局部热源均从相同的单位端口复场叠加获得，
不另外训练不相关的独立模型。
材料界面使用分域基和相应跳跃条件，
不得以单个光滑神经输出跨越全部材料界面。

\subsection{几何等变性与外场}
当背景各向同性且没有固定世界方向输入，
整体旋转和平移时矢量场应协同变换，
端口标量响应保持不变；但固定海流、重力、
各向异性电导率张量和世界外场需同步变换才能保留该对称性。
局部装配件任意改变相对姿态不是整体不变变换：
模型必须真实计算新的耦合关系。

\section{新算法：BFZI 条件化保结构低阶响应算子}\label{sec:neural-rom}
\subsection{目标：学习解空间，而不是学习互不相容的端口数值}
第 \ref{sec:galerkin} 节构造了严格同源的海水电流系数、
复阻抗和空间焦耳热。本节进一步定义一种\textbf{可训练但结构可证明}的求解路线。
在给定合法场景 \(\mathcal G\) 中，先构造一个较丰富的
几何适配、连续、散度自由和界面法向零通量的海水电流字典：
\[
\bm W_0(\bm x;\mathcal G)
=[\bm w_1(\bm x;\mathcal G),\ldots,\bm w_M(\bm x;\mathcal G)],
\qquad \bm w_m\in\mathcal V(\Omega_s).
\]
一个候选构造是局部矢量势 curl 模态、曲面切向边界适配模态、
低阶全局环流模态的组合；\textbf{在任意六自由度非凸封装/多连通海水域下
如何高效构造此字典，是尚未解决的算法问题，不得假称已通用实现}。
只使用带环境方向信息的物理合法局部坐标映射，
避免将固体与海水的边界条件用同一个不分区网络处理。

计算字典级真实物理 Gram 矩阵
\[
\begin{aligned}
[\bm R_0]_{mn}&=\int_{\Omega_s}\sigma_s^{-1}\bm w_m\cdot\bm w_n\,\mathrm dV,\\
[\bm L_0]_{mn}&=\mu_0\int_{\Omega_s}\bm w_m\cdot\mathcal G_0\bm w_n\,\mathrm dV,\\
[\bm K_0]_{mp}&=\mu_0\int_{\Omega_s}\bm w_m\cdot\bm C_p^{\rm BFZI}\,\mathrm dV.
\end{aligned}
\]
其中 \(\bm R_0\succ0,\bm L_0\succeq0\) 的充分前提为
实值线性独立字典、正电导率和适用的磁核。

令一个轻量条件网络接收
\(c=(\bm g_{\rm BFZI},\bm g_{\rm pack},q_{\rm rel},
\sigma_s,\omega,\text{thermal descriptors})\)，输出
\begin{equation}
\boxed{\bm U_\theta(c)\in\mathbb R^{M\times r},\qquad r\ll M.}
\label{eq:learned-subspace}
\end{equation}
为防止列空间退化，可将 \(\bm U_\theta\)
在 \(\bm R_0\) 内积下正交化，令
\[
\bm S_\theta=
\bm U_\theta
\bigl(\bm U_\theta^{\mathsf T}\bm R_0\bm U_\theta\bigr)^{-1/2},
\quad
\bm S_\theta^{\mathsf T}\bm R_0\bm S_\theta=\mathsf I,
\]
必要时用正交 QR/SVD 及秩容差稳定实现。
定义低维连续电流基
\begin{equation}
\boxed{\bm W_\theta(\bm x;\mathcal G,c)
=\bm W_0(\bm x;\mathcal G)\bm S_\theta(c).}
\label{eq:learned-current-basis}
\end{equation}
由于 \(\bm S_\theta\) 对空间 \(\bm x\) 是常数，
\(\bm W_\theta\) 仍满足字典空间的散度和界面法向约束。
\textbf{神经网络只选择/混合已满足物理约束的电流空间，
不直接写出任意的 \(R,L,M\) 或电流。}

\subsection{快速矩阵收缩与端口/场/热源同源解码}
在字典级矩阵存在时，低维组装只需小矩阵收缩
\begin{equation}
\boxed{
\begin{aligned}
\bm R_\theta&=\bm S_\theta^{\mathsf T}\bm R_0\bm S_\theta,\\
\bm L_\theta&=\bm S_\theta^{\mathsf T}\bm L_0\bm S_\theta,\\
\bm K_\theta&=\bm S_\theta^{\mathsf T}\bm K_0.
\end{aligned}}\label{eq:structured-contractions}
\end{equation}
因此
\[
\bm A_\theta=\bm R_\theta+\jj\omega\bm L_\theta,\qquad
\bm z_\theta=-\jj\omega\bm A_\theta^{-1}\bm K_\theta\bm I,\qquad
\widehat{\bm J}_s(\bm x)=\bm W_0(\bm x)\bm S_\theta\bm z_\theta,
\]
并通过式 \eqref{eq:zsw-schur} 得到
\(\widehat{\bm Z}_{\rm sw}\)、通过式 \eqref{eq:sea-power-proof}
得到 \(\widehat q_s(\bm x)\)。
即使 \(\bm S_\theta\) 来自神经网络，
只要所有矩阵和场读出用同一组字典及同一 \(\bm S_\theta\)，
第 \ref{sec:galerkin} 节的互易、被动、功率等式\textbf{仍然成立}。
这是比额外加入若干物理残差惩罚更强的结构性结果。

\textbf{温度相关性}：若 \(\sigma_s=\sigma_s(T)\)，
必须更新 \(\bm R_0(T)\) 及其正交归一，
不能把参考温度的 \(\bm R_0\) 缓存跨所有温度直接使用。
如果字典位置随装配件运动改变，
\(\bm R_0,\bm L_0,\bm K_0\) 同样可能需要重组。
可采用训练覆盖域内的参数化核插值、缩放律与自适应缓存，
但插值/近似后的结构约束与真实空间积分是否一致需另行处理。

\subsection{BFZI 结构如何加速 \(\bm K_0\)}
用 BFZI 主螺旋匝号仿射关系，
在固定节点组 \(b,\mu\) 记
\(\bm r_{p,b,k}=\bm r_{p,b,0}+k\Delta\bm r_{p,b}\)，
\(\bm h_{p,b,k}=\bm h_{p,b,0}+k\Delta\bm h_{p,b}\)，
则
\begin{equation}
[\bm K_0]_{mp}
=\frac{\mu_0}{4\pi}
\int_{\Omega_s}\bm w_m(\bm x)\cdot
\left[
\sum_{b,\mu}\sum_{k=0}^{N_p-1}
\frac{\bm h_{p,b,k}}{\|\bm x-\bm r_{p,b,k}\|}
+\bm C_p^{\rm dir}(\bm x)\,4\pi
\right]\mathrm dV.
\label{eq:bfzi-K-assembly}
\end{equation}
括号内第一项是主螺旋的\textbf{单匝号核聚合}；
可复用 BFZI 论文的精确特殊情形与有余项界的分块法。
\(\bm C_p^{\rm dir}\) 是非重复馈线/回流对原 \(\bm C_p\)
的贡献，故第二项中的 \(4\pi\) 用于统一第一项记号。
可以把字典模式与 BFZI 端口线源的积分矩预先分解/缓存，
对空间点进行批量查询。

\textbf{限制}：BFZI 匝号加速只减少\textbf{主螺旋源项/线元配对}计算，
不自动消除 \(\bm R_0,\bm L_0\) 的体积/边界积分代价，
也不能保证任意几何下总复杂度与匝数无关。
真正的整体加速须单独测量字典构建、积分矩组装、
神经子空间生成、\(r\times r\) 求解与任意空间点查询。

\subsection{有约束训练和有限维扰动控制}
可用高精度教师的单位端口涡流场 \(\bm J_p^\star\)
训练 \(\bm S_\theta\)，令其直接最小化
\[
\mathcal L_{\rm state}
=\sum_p\int_{\Omega_s}
\frac{\|\widehat{\bm J}_{s,p}-\bm J_p^\star\|^2}
{\sigma_s}\,\mathrm dV,
\]
并同时使用端口复阻抗和热点场数据监督；
不需要额外训练一个互不一致的阻抗网络。
若字典空间的矩阵组装存在可控近似
\(\widehat{\bm A}=\bm A+\Delta\bm A\),
\(\widehat{\bm K}=\bm K+\Delta\bm K\)
且矩阵尺寸、基空间、基向量意义均一致，
则由预解式恒等式
\(\widehat{\bm A}^{-1}-\bm A^{-1}
=-\bm A^{-1}\Delta\bm A\,\widehat{\bm A}^{-1}\)，
有
\begin{align}
\|\widehat{\bm Z}_{\rm sw}-\bm Z_{\rm sw}\|
\leq \omega^2\bigl(&
\|\Delta\bm K\|\|\widehat{\bm A}^{-1}\|\|\widehat{\bm K}\|\nonumber\\
&+\|\bm K\|\|\bm A^{-1}\|\|\Delta\bm A\|
\|\widehat{\bm A}^{-1}\|\|\widehat{\bm K}\|\nonumber\\
&+\|\bm K\|\|\bm A^{-1}\|\|\Delta\bm K\|
\bigr).\label{eq:z-perturbation}
\end{align}
这是\textbf{同一有限维空间内的组装扰动界}，
不是跨不同几何、不同基空间时的直接误差保证；
投影截断误差仍须独立估计。

\subsection{两种运行路径与物理一致性失效点}
\textbf{REFERENCE}：原始 BFZI 几何、完整材料与边界模型、
足够丰富的合规电流字典/高精度空间解；
\textbf{FAST}：网络选择子空间并解低维被动系统，适当缓存积分核。
在 FAST 中也可进一步让网络直接预测结构化低维
\(\bm R,\bm L,\bm K\) 加速，但若 \(\bm R\) 不再等于所解码
\(\bm J_s\) 的焦耳内积，式 \eqref{eq:sea-power-proof}
就不再成立；仅约束正定\textbf{不等于空间热源一致}。
因此这一捷径不能直接用于核心方法，除非额外设计
有局部正半定损耗分解且与场解码一致的严格闭合机制。

\section{保半正定的核装配与 BFZI 在线计算}\label{sec:positive-assembly}
\subsection{瓶颈和目标}
第 \ref{sec:admissible-atlas} 节的局部图谱减少了
跨几何重建 FEM 网格的需求，但 \(\bm R,\bm L,\bm K\) 的空间积分
仍可能成为推理瓶颈。因此需要在\textbf{保持物理 Gram 结构}的前提下
改造核装配。下面区分两种来源：
\textbf{精确对称积分 / 一致离散}和
\textbf{用正权核特征近似后仍具有半正定结构}。
不能通过独立预测矩阵元素来保证整体功率守恒。

\subsection{电阻 Gram：正权局部求积}
对同一套几何条件化实电流基 \(\bm W_i\)，
在实际导电域选取正权求积
\((\bm x_\nu,w_\nu)\)，\(w_\nu>0\)，并定义
\begin{equation}
\widehat{\bm R}_{ij}
=\sum_{\nu=1}^{Q_R}
\frac{w_\nu}{\sigma_s(\bm x_\nu,T)}
\bm W_i(\bm x_\nu)\cdot\bm W_j(\bm x_\nu).
\label{eq:R-positive}
\end{equation}
若用 \(\bm D_R\) 表示把每个模式在所有求积点的三分量
叠加后形成的加权采样矩阵，则
\[
\widehat{\bm R}=\bm D_R^{\mathsf T}\bm D_R\succeq0.
\]
若在采样点上字典列仍线性独立，则 \(\widehat{\bm R}\succ0\)；
若不独立，应消除冗余模态、细化近界面求积或选择稳定秩，
而不是在未解释空间损耗的情况下随意添加对角电阻。
用\textbf{同一求积节点}的
\(q_{s,\nu}=|\sum_i z_i\bm W_i(\bm x_\nu)|^2/(2\sigma_s)\)，
可保证
\[
\tfrac12\bm z^{\mathrm H}\widehat{\bm R}\bm z
=\sum_\nu w_\nu q_{s,\nu}.
\]
这只是\textbf{离散求积功率等式}。
如果用其它积分规则或连续精确热积分，
误差由 \(\widehat{\bm R}-\bm R\) 控制，
不得把离散等式当成零误差的连续物理结果。

\subsection{磁核 Gram：共同的非负 Fourier 特征}
在均匀自由空间磁导率、场足够衰减且基函数可作 Fourier 变换时，
记 \(\widehat{\bm W}_m(\bm k)=
\int_{\mathbb R^3}\bm W_m(\bm x)e^{-\jj\bm k\cdot\bm x}\,\mathrm d\bm x\)。
对 Coulomb Green 核
\(G_0(\bm x)=1/(4\pi\|\bm x\|)\)
有形式上
\[
\widehat{G_0}(\bm k)=\|\bm k\|^{-2}.
\]
在满足积分收敛的条件下，
由 Parseval 恒等式
\begin{equation}
[\bm L]_{mn}
=\frac{\mu_0}{(2\pi)^3}
\int_{\mathbb R^3}
\frac{\operatorname{Re}[
\widehat{\bm W}_m(\bm k)^*\cdot
\widehat{\bm W}_n(\bm k)]}
{\|\bm k\|^2}\,\mathrm d\bm k.
\label{eq:magnetic-fourier}
\end{equation}
该表达在 \(\bm k=0\) 附近的可积性、
高频收敛和外海水模式必须由实际函数空间保证。
用排除零波数奇点并匹配正权的 Fourier 节点
\((\bm k_\nu,a_\nu)\)、\(a_\nu>0\)，定义
\begin{equation}
\boxed{
\widehat{\bm L}_{mn}=
\mu_0\sum_{\nu=1}^{Q_L}
\frac{a_\nu}{\|\bm k_\nu\|^2}
\operatorname{Re}\!\left[
\widehat{\bm W}_m(\bm k_\nu)^*\cdot
\widehat{\bm W}_n(\bm k_\nu)\right].
}
\label{eq:L-positive}
\end{equation}
对于所有实系数向量 \(\bm c\)，
\[
\bm c^{\mathsf T}\widehat{\bm L}\bm c
=\mu_0\sum_\nu
\frac{a_\nu}{\|\bm k_\nu\|^2}
\left\|\sum_m c_m\widehat{\bm W}_m(\bm k_\nu)\right\|^2
\ge0.
\]
因此只要所有模式使用\textbf{同一组正权谱特征}，
即使求积尚有误差，所得实对称磁矩阵仍半正定。
实际实现可用带 \(\pm\bm k_\nu\) 成对的实/虚特征
保证矩阵实对称；\textbf{此构造是标准正定核性质的具体应用，
不是首次发现 Fourier 保正性的数学定理。}

对铜载流提升场、盒壁模式和海水模式，
可把全部电流函数列组成 \(\bm V=[\bm U,\bm W]\)，
用完全相同的 Fourier 特征生成完整
\(\widehat{\bm L}_{ab}\) 块；
这样\textbf{全部块联合半正定}，
联合 Schur 补的被动性和离散功率证明继续适用。
若海水 \(\bm K\) 使用精确 BFZI 线积分、
\(\bm L_{00}\) 使用另一组独立近似，
则不能自动断言合并后的完整磁 Gram 为半正定；
必须核对共同能量内积或回到联合同源组装。

\subsection{刚体变换可复用的谱特征}
对随组件刚体移动的模式
\(\bm W_p(\bm x)
=R_p\widehat{\bm W}_p(R_p^{\mathsf T}(\bm x-\bm c_p))\)，
其 Fourier 变换严格满足
\begin{equation}
\boxed{
\widehat{\bm W}_p(\bm k)
=e^{-\jj\bm k\cdot\bm c_p}\,
R_p\widehat{\widehat{\bm W}}_p(R_p^{\mathsf T}\bm k).
}
\label{eq:rigid-fourier}
\end{equation}
这里 \(\widehat{\widehat{\bm W}}\) 表示\textbf{局部参考模式的 Fourier 变换}，
不表示再对 Fourier 数据作一次变换。
因此局部几何固定且其模式不侵犯另一实体时，
只需改变相位、方向和频点采样，就能更新相对位姿的谱 Gram。
在任意旋转下，参考谱需要在旋转后的
\(R_p^{\mathsf T}\bm k_\nu\) 处评估；
插值误差可改变数值精度，
但若最终 \(\widehat{\bm L}\) 仍由同一组采样特征
形成 Gram，半正定性仍得到保持。
相对位姿变化不会使跨组件磁耦合不变；
\(\bm c_2-\bm c_1\) 对应的相位仍需计算。

\subsection{BFZI 快速构造端口耦合 \(\bm K\)}
第 \ref{sec:galerkin} 节的
\(\bm K_{mp}=\mu_0\int_{\Omega_s}\bm W_m\cdot\bm C_p\,dV\)
由 Green 核互易性可改写为
\begin{equation}
\boxed{
K_{mp}=\int_{\Gamma_p}\bm A_m^{\rm mode}(\bm r_p)
\cdot \mathrm d\bm\ell_p,\quad
\bm A_m^{\rm mode}(\bm r)
=\mu_0\int_{\Omega_s}G_0(\bm r-\bm x)\bm W_m(\bm x)\,\mathrm dV.
}
\label{eq:K-reciprocal}
\end{equation}
这表明可以\textbf{预计算或快速查询电流模式的矢量势}，
再沿 BFZI 连续线圈执行结构化线积分，
不必在每次端口查询中完整求一个海水三重体积分。
在固定线段组与 Gauss 节点，BFZI 仍满足
\(\bm r_{b,k}=\bm r_{b,0}+k\Delta\bm r_b\) 和
\(\bm h_{b,k}=\bm h_{b,0}+k\Delta\bm h_b\)。
若 \(\bm A_m^{\rm mode}\) 在匝号块覆盖的路径邻域
具有可控导数，可按 \(k\) 做 Taylor / Chebyshev 模式求和，
用
\(\sum_{k\in T}(k-k_c)^j\) 的显式矩求和加速；
在近场陡变、靠近盒边和模式奇异位置回退直接求积。
注意这与原 BFZI 论文对 \(1/\|\bm x-\bm r_k\|\) 核做二阶块求和
是\textbf{扩展应用}，新的 \(\bm A_m^{\rm mode}\) 误差界需单独推导，
不能沿用原文的余项常数。

\subsection{正定性、连续损耗与加速之间的取舍}
\textbf{路径 A（首选保结构）}：
所有 \(\bm R\) 均来自实际分区上的正权积分，
所有 \(\bm L\) 块来自同一正权谱/空间核 Gram；
因此任意子空间选择与联合 Schur 补保持离散被动性，
输出空间 \(q_i(\bm x)\) 由同一当前电流模式求得。
\textbf{路径 B（严格源核）}：
BFZI 互易线积分快速给出精确/有界的 \(\bm K\)。
此时对只有海水的低维散射阻抗，独立
\(\bm R_s\succ0,\bm L_s\succeq0\) 足以保证
海水散射部分被动；但若要把铜/盒壁共同耦合进完整 Schur，
则必须继续检查\textbf{完整磁 Gram 的一致性}。
\textbf{禁用的捷径}：
单独神经预测 PSD \(\bm R\) 和各点损耗却不给出相同的场内积，
不能宣称端口积分功率严格相等。

\subsection{在线复杂度的正确统计}
若字典模式数 \(M\)、学习子空间维数 \(r\)、电阻求积点数
\(Q_R\)、谱节点数 \(Q_L\)，且
\(\bm D_R,\bm D_L\) 可在当前场景生成，
则可先计算小型特征
\[
\bm F_{R,r}=\bm D_R\bm S_\theta,\qquad
\bm F_{L,r}=\bm D_L\bm S_\theta,
\quad
\bm R_r=\bm F_{R,r}^{\mathsf T}\bm F_{R,r},\quad
\bm L_r=\bm F_{L,r}^{\mathsf T}\bm F_{L,r},
\]
从而避免显式构造 \(M\times M\) 全矩阵。
典型稠密运算的贡献为
\(O((Q_R+Q_L)Mr+(Q_R+Q_L)r^2+r^3)\)，
\textbf{不包含}字典建立、特征值评估、
BFZI \(\bm K\) 核装配及最终场查询。
若分区紧支撑，\(\bm D_R\) 实际稀疏，
应利用局部结构减少常数和矩阵乘法代价。
要获得真实在线加速，必须使
\(Q_R,Q_L,M\) 与图谱计算在目标参数域内有可复用/分离结构；
仅令 \(r\) 很小并不足够。

\subsection{几何重组时的失败回退}
相对位姿导致局部模式支撑侵犯另一盒体、
间隙接近零、正权采样不足、低阶电流基不满秩或
谱/局部核逼近无法达到容差时：
增广间隙图谱、加密局部求积、缩小模型适用域
或回退高精度场求解。
不能为维持速度而输出不合规的电流或负功率损耗。

\section{电磁低阶误差界、几何回退与训练最小闭环}\label{sec:error-algorithm}
\subsection{Galerkin 电流误差的能量范数}
固定某个合法的 BFZI--封装--海水场景 \(\mathcal G\)，
在一个有限维合规的高精度参考电流字典上，定义
\(\bm A=\bm R+\jj\omega\bm L\)，
\(\bm R\succ0\)、\(\bm L\succeq0\)，
\(\bm A\bm z^\star=\bm f\)；
设由神经选择空间投影得到的近似解为 \(\bm z_r\)
（已提升到同一参考字典坐标），
\(\bm r_{\rm full}=\bm f-\bm A\bm z_r\)，
\(\bm e=\bm z^\star-\bm z_r\)。
用 \(\|\bm u\|_{\bm R}
=(\bm u^{\mathrm H}\bm R\bm u)^{1/2}\)
和 \(\|\bm v\|_{\bm R^{-1}}\) 表示能量范数及其对偶范数。
由于 \(\bm A\bm e=\bm r_{\rm full}\)，有
\[
\|\bm e\|_{\bm R}^2
=\operatorname{Re}(\bm e^{\mathrm H}\bm A\bm e)
=\operatorname{Re}(\bm e^{\mathrm H}\bm r_{\rm full})
\le \|\bm e\|_{\bm R}\|\bm r_{\rm full}\|_{\bm R^{-1}},
\]
所以
\begin{equation}
\boxed{
\|\bm z^\star-\bm z_r\|_{\bm R}
\le \|\bm r_{\rm full}\|_{\bm R^{-1}}.
}
\label{eq:galerkin-residual-bound}
\end{equation}
这是一条\textbf{同一高精度有限维参考空间}下的
标准残差能量误差界，适用于任何合规的近似电流状态，
不要求神经网络满足特殊训练损失。
其代价是完整残差和对偶范数可能在线昂贵；
可在代表性教师场景/抽样几何中精确计算用于验收，
在线使用参数化低维残差估计与阈值回退。
若 \(\bm A,\bm f\) 自身是近似装配，
必须将装配误差纳入 \(\bm r_{\rm full}\) 或另外增加误差项，
不能把式 \eqref{eq:galerkin-residual-bound}
直接应用于未测量的真实连续算子。

\subsection{最优子空间近似与条件数}
对于带 \(\bm R\) 内积的同一字典，
定义 \(\bm B=\bm R^{-1/2}\bm L\bm R^{-1/2}\)，
其为实对称半正定。
\[
\bm R^{-1/2}\bm A\bm R^{-1/2}
=\mathsf I+\jj\omega\bm B.
\]
该变换算子在能量范数内的算子范数为
\(\sqrt{1+\omega^2\|\bm B\|_2^2}\)，
实部强制性常数为 \(1\)，
因而标准复数 Galerkin 估计提供
\begin{equation}
\|\bm e\|_{\bm R}
\le
\sqrt{1+\omega^2\|\bm B\|_2^2}
\inf_{\bm v\in\mathcal V_r}
\|\bm z^\star-\bm v\|_{\bm R}.
\label{eq:cea-complex}
\end{equation}
这是较保守的上界，
仅用于分析\textbf{频率放大效应和子空间覆盖不足}，
不能据此预测实际误差总能保持在某固定百分比内。
在高频、海水电导率高或皮肤深度显著缩小时，
\(\omega\|\bm B\|\) 可能增大，应相应增广涡流模式与训练频段。

\subsection{端口/热源的误差与残差}
对固定的同一高精度字典及精确耦合 \(\bm K\)，
低阶海水阻抗误差来源于响应差
\(\Delta\bm z\)：
\(\Delta\bm v_s=\jj\omega\bm K^{\mathsf T}\Delta\bm z\)，
故
\[
\|\Delta\bm v_s\|_2
\le |\omega|\,\|\bm K^{\mathsf T}\bm R^{-1/2}\|_2
\|\Delta\bm z\|_{\bm R}.
\]
这给出从同源电流的能量误差到端口电压误差的界。
对于同一欧姆积分，海水功率误差满足
\begin{equation}
\left|
\tfrac12\|\bm z^\star\|_{\bm R}^2
-\tfrac12\|\bm z_r\|_{\bm R}^2
\right|
\le \tfrac12
(\|\bm z^\star\|_{\bm R}+\|\bm z_r\|_{\bm R})
\|\bm z^\star-\bm z_r\|_{\bm R}.
\label{eq:loss-error}
\end{equation}
但\textbf{最大局部热源、盒壁热点和温度峰值}
一般不能仅由一个全局 \(L^2\) 电流误差界控制；
仍需局部增强采样、近界面场验证和热传播稳定性分析。
不同几何下的高精度字典不一致时，
需先指定稳定的场映射/投影，才能比较其范数。

\subsection{一条可计算的最小运行流程}
在物理参考求解器可用后，按下列算法推进；
这里的 \texttt{REFERENCE} 与 \texttt{FAST}
仍是严格分开的模型层级，不能互相偷偷代替。

\paragraph{离线阶段}
确定 BFZI 几何合法域和导体/馈回流约定；
为典型封装和海水构造合规局部 curl/Piola 图谱及拓扑模式；
生成参考场景、端口单位电流教师样本；
计算自洽参考场、材料损耗和 \(\bm R,\bm L,\bm K\)；
将几何--频率--位姿--材料条件映射到可表达该响应的低维子空间，
训练 \(\bm S_\theta\) 选择器，并验证
式 \eqref{eq:galerkin-residual-bound} 对应的电流误差与空间误差。
对于几何改变后不能复用的矩阵核，记录在线必须重组的具体块。

\paragraph{在线阶段}
\begin{enumerate}
\item 读取 BFZI 线圈和各组件的完整尺寸、材料、相对六自由度位姿；
检查非法穿透、几何非退化、局部图谱的正间隙条件。
\item 选择近盒、间隙、远场和必要拓扑电流模式；
将每个模式映射至真实世界坐标并保留连续性/法向零通量。
\item 用正权物理积分或共同的磁核谱特征
组装所需的电阻、磁能和 BFZI 端口耦合矩阵；
主螺旋源项执行有误差控制的匝号聚合。
\item 用网络选择低阶子空间 \(\bm S_\theta\)，
进行小矩阵收缩及 \(r\) 维复线性系统求解；
若含有限铜导体/导电盒壁则使用联合 Schur 补，
不可独立叠加修正。
\item 由同一低维电流状态计算复阻抗、连续空间
\(\bm J,\bm E,\bm B\) 和逐材料 \(q_i\)；
向电路与守恒热 ODE 传入同一组损耗投影。
\item 若场景超出图谱覆盖、矩阵失去秩、
边界近接触或样本内残差/局部场误差失控，
增广模式或切换参考求解。
\end{enumerate}
在线查询远离线圈与海水模式支持集的
\(\bm E,\bm B\) 仍需要连续 Green 核/多极或谱快速场评价；
\textbf{小维度求出涡流系数不等于所有三维场查询都已经快速。}

\subsection{验证的优先级}
首个可复现实例限定：两只 BFZI 合法线圈、
非导电盒壁、均匀静止海水、一个频率、
已知六自由度相对位姿；
不先训练大网络。
先证明/对照：散度与法向通量、
互易、端口阻性功率与海水损耗的离散恒等式、
参考场的空间误差和实际组装耗时。
随后分别增加频率扫、间隙趋小、封装尺寸变化、
有限铜截面、导电盒壁和热时间演化。
研究记录优先关注\textbf{物理计算是否正确、线上矩阵组装是否有真实效率优势}，
无需添加大量额外审计程序。

\section{材料损耗、互易、被动性与功率一致}\label{sec:loss}
\subsection{完整端口矩阵和材料损耗归属}
设 \(\bm Z_{\rm EM}\) 是从完整电磁参考面看到的
无源线性耦合器内部端口矩阵，\textbf{不含外部负载和补偿网络}。
铜导体、导电海水、盒壁、屏蔽及磁芯耗散分别由不重叠材料域负责；
相量峰值约定下
\begin{equation}
q_i(\bm x)=\frac12\operatorname{Re}
[\bm J_i(\bm x)\cdot\bm E_i^*(\bm x)]
\quad\text{(欧姆耗散区域)}.
\end{equation}
频散或磁滞材料须使用适当周期平均损耗模型，
不可将同一损耗同时计入磁化功和电流欧姆项。
电磁总耗散应满足
\begin{equation}
\boxed{
P_{\rm EM,int}=\frac12\bm I^{\mathrm H}
\Herm{\bm Z_{\rm EM}}\bm I
=\sum_i\int_{\Omega_i}q_i(\bm x)\,\mathrm dV.
}\label{eq:power-balance}
\end{equation}
该等式成立的条件是两侧采用相同电流/端口参考和
完整且不重复的内部耗散分区。
不能用原 BFZI 海水反作用阻抗实部代表包括铜损在内的总损耗。

\subsection{结构约束}
在线性互易材料条件下要求 \(\bm Z=\bm Z^{\mathsf T}\)；
非互易介质下不强制对称。
无源端口要求
\[
\Herm{\bm Z_{\rm EM}}\succeq0.
\]
注意互阻实部可能为负数，不能逐元素裁剪电阻。
若以 \(\bm L\bm L^{\mathrm H}\) 参数化耗散矩阵，
还应使其与实际空间材料损耗和端口电流相容。
宽带应用必须进一步满足因果性和材料频散限制；
单频正半定不足以证明完整频带的物理正确性。

\subsection{监督与物理约束}
训练损失可同时包含端口数据、场数据、空间损耗、
电流连续/边界条件、\eqref{eq:power-balance} 的功率残差。
优先通过解码结构保证强物理约束，
而不是叠加大量数值惩罚掩盖物理不一致。

\section{新推导：导体、海水与导电盒壁的统一端口 Schur 补}\label{sec:multidomain-schur}
\subsection{为何不能直接相加独立损耗修正}
只对理想丝状线圈与线性均匀海水，上一节海水
\(\bm Z_{\rm sw}\) 可作为空气基底上的附加散射项。
若有限截面导体、金属屏蔽和海水同时存在，
它们的感应电流会通过磁场相互作用；将各区域
\emph{独立求解}的 \(\Delta\bm Z_{\rm Cu},\Delta\bm Z_{\rm box},
\Delta\bm Z_{\rm sea}\) 逐项相加一般\textbf{不是}完整耦合解。
应使用同一低阶块矩阵描述全部导电区域。

\subsection{端口载流提升场与无端口环流场}
令 \(\Omega_c=\bigcup_i\Omega_i^{\rm cond}\)
为所有线性导电实体（包括铜线、海水和可选的导电盒壁）
的互不重叠并集。导体端口参考面、端点与回路电流方向必须明确。
构造单位端口载流提升场
\(\bm U_p(\bm x)\)，具有正确的第 \(p\) 个端口通量和绝缘面
法向零电流，满足实体内部电流连续性；
同时构造零端口电流的散度自由纠正基
\(\bm W_m(\bm x)\)。
写作
\begin{equation}
\boxed{
\bm J(\bm x)
=\bm U(\bm x)\bm I+\bm W(\bm x)\bm z.
}
\label{eq:joint-current}
\end{equation}
这里 \(\bm U\) 的端口通量是一种约束仿射提升，
而 \(\bm W\) 的端口净通量严格为零。
对海水等无外接导电端口的区域，
\(\bm U\) 在该材料域可以为零但 \(\bm W\) 可以非零。
实际固体的接触导电性、连接关系、绝缘涂层和
海水--金属界面需要按真实物理设置；不能假设全部材料都电绝缘。

\subsection{统一阻性与磁性双线性型}
设 \(\bm V_0(\bm x)=\bm U(\bm x)\)、
\(\bm V_1(\bm x)=\bm W(\bm x)\)，
对两组基函数 \(a,b\in\{0,1\}\) 定义
\begin{align}
[\bm R_{ab}]_{mn}
&=\int_{\Omega_c}\sigma^{-1}
\bm V_{a,m}\cdot\bm V_{b,n}\,\mathrm dV,
\\
[\bm L_{ab}]_{mn}
&=\mu_0\int_{\Omega_c}
\bm V_{a,m}\cdot\mathcal G_0\bm V_{b,n}\,\mathrm dV.
\end{align}
记 \(\bm A_{ab}=\bm R_{ab}+\jj\omega\bm L_{ab}\)。
对合规的分段材料空间和无源实导电率，
\(\bm R_{ab}=\bm R_{ba}^{\mathsf T}\),
\(\bm L_{ab}=\bm L_{ba}^{\mathsf T}\)，
\(\bm R_{11}\succ0\)（对线性独立纠正基）。
由带端口约束的谐波弱式获得
\begin{equation}
\begin{bmatrix}
\bm v\\ \bm 0
\end{bmatrix}
=
\begin{bmatrix}
\bm A_{00}&\bm A_{01}\\
\bm A_{10}&\bm A_{11}
\end{bmatrix}
\begin{bmatrix}
\bm I\\\bm z
\end{bmatrix},\qquad
\bm z=-\bm A_{11}^{-1}\bm A_{10}\bm I .
\label{eq:joint-system}
\end{equation}
矩阵第二行表示对全部\emph{零端口电流测试场}的电磁方程残差为零。
对端口载流提升 \(\bm U\) 则由第一行读出端口电压；
严格性依赖完整端口边界及源项/电势约定。

消去 \(\bm z\) 后得到
\begin{equation}
\boxed{
\bm Z_{\rm cond,total}
=\bm A_{00}-\bm A_{01}\bm A_{11}^{-1}\bm A_{10}.
}
\label{eq:joint-schur}
\end{equation}
这包含导体皮肤/邻近重分布、海水涡流和盒壁涡流
在单一线性 MQS 模型中的\textbf{耦合}影响。
不能再把它与已经由同一模型包含的海水/铜损阻抗重复叠加。
当载流提升退化成理想无耗散丝状电流、
\(\bm R_{00}=\bm R_{01}=0\)，
\(\bm A_{00}=\jj\omega\bm L^{\rm air}\)，
\(\bm A_{10}=\jj\omega\bm K\) 时，
式 \eqref{eq:joint-schur} 恰退化为
\(\jj\omega\bm L^{\rm air}+
\omega^2\bm K^{\mathsf T}
(\bm R_s+\jj\omega\bm L_s)^{-1}\bm K\)。

\subsection{被动性、互易与任意导电材料的功率一致性}
上述块矩阵复对称，且 Schur 补保持转置对称：
\(\bm Z_{\rm cond,total}^{\mathsf T}=\bm Z_{\rm cond,total}\)。
对实际解 \(\bm y=[\bm I;\bm z]\)，
用式 \eqref{eq:joint-system} 得
\[
\bm I^{\mathrm H}\bm v
=\bm y^{\mathrm H}
(\bm R+\jj\omega\bm L)\bm y,
\quad
\bm R=
\begin{bmatrix}\bm R_{00}&\bm R_{01}\\
\bm R_{10}&\bm R_{11}\end{bmatrix}.
\]
取实部得到
\begin{equation}
\boxed{
\frac12\bm I^{\mathrm H}
\operatorname{Herm}(\bm Z_{\rm cond,total})\bm I
=\frac12\bm y^{\mathrm H}\bm R\bm y
=\sum_i\frac12\int_{\Omega_i^{\rm cond}}
\frac{|\bm J_i(\bm x)|^2}{\sigma_i(\bm x)}\,\mathrm dV\ge0.
}
\label{eq:joint-power}
\end{equation}
因此完整端口矩阵被动、互易（在本节互易条件下），
而且按真实实体域划分的焦耳热\textbf{不会相互重复计入}。
此证明比人为调低多种独立损耗预测的误差更有结构保证。

\subsection{磁性材料、频散及非导电介质的边界}
本节的磁核使用均匀 \(\mu_0\)；
磁芯非线性、磁滞、频散、磁化电流及位移电流
\textbf{不由单纯的 \(\bm R,\bm L\) 块自动覆盖}。
在磁性部件存在时，应扩展有限维状态到磁化/磁通变量，
再推导相应被动储能--耗散形式与额外材料热源。
只在该扩展模型验证后才可宣称适用任意磁芯。
非导电灌封本身无欧姆体积损耗，但影响海水通道与热传导；
其作用必须保留在几何/边界条件而非虚构电阻项中。

\subsection{与 BFZI 几何和热演化的连接}
\(\bm U_p\) 的主载流分量继承 BFZI 直边--角部
连续路径及匝号仿射源项，并显式处理有限导体截面及馈/回流；
用 BFZI 结构化核积分为块矩阵
\(\bm A_{00},\bm A_{01}\) 的磁场项提供加速。
\(\bm W\) 的低维空间可由第 \ref{sec:neural-rom} 节的
几何条件化模型选择。
同一 \(\bm y\) 计算各区
\(q_i(\bm x)=|\bm J_i(\bm x)|^2/(2\sigma_i)\)，
再将 \(q_i\) 投影到式 \eqref{eq:thermal-rom} 的热状态输入，
从而明确组成电磁--电路--热的闭环动力系统。

\section{海水耦合瞬态热场与守恒降阶演化}\label{sec:thermal}
\subsection{固体与海水的热方程}
对封装、导体、磁芯等固体，在自身材料坐标内
\begin{equation}
\rho_i c_i(T)\,\mathrm D_t T_i
=\grad\cdot(k_i(T)\grad T_i)+q_i^{\rm EM}.
\label{eq:heat-solid}
\end{equation}
对于海水
\begin{equation}
\rho_sc_s(T_s)[\partial_tT_s+\bm v_s\cdot\grad T_s]
=\grad\cdot(k_s(T_s)\grad T_s)+q_s^{\rm EM}.
\label{eq:heat-sea}
\end{equation}
界面给出接触温差/接触热阻及热通量平衡；
海水远场具有背景温度/热输运条件。
第一阶段可采用静止海水或指定流速场，
但不能宣称不建流场就能正确预测任意湍流与自然对流。
运动组件中的温度采用材料导数，避免将刚体运输重复计入热源。

\subsection{几何条件化连续温度场}
对第 \(i\) 个组件使用随体坐标
\(\bm\xi_i=R_i^{\mathsf T}(\bm x-\bm c_i)\)，
研究连续低维表达
\begin{equation}
T_i(\bm x,t)=T_{\rm ref}(\bm x,t)
+\sum_{\ell=1}^{r_i}
a_{i,\ell}(t)\Phi_{i,\ell}(\bm\xi_i;\bm g_i).
\label{eq:thermal-basis}
\end{equation}
\(\Phi\) 可来自解析 Green/传热本征基与有限维学习修正，
需对边界接口和热容量保持物理一致。
\textbf{不可默认固定 FEM 网格 POD 基支持全部 BFZI 形状与盒体尺寸}。
海水的热尾迹应使用独立的连续核/低阶状态表达，
否则可能错误丢失流向下游的热能。

\subsection{守恒低阶热动力系统}
对热量方程进行几何条件化投影，得到候选
\begin{equation}
\boxed{
\bm M(\bm g,\bm a)\dot{\bm a}
+\bm K(\bm g,\bm a,\bm q,\bm v_s)\bm a
=\bm B(\bm g)\bm p_{\rm EM}
+\bm f_\theta(\bm a,\bm g,\bm q,\bm v_s).
}
\label{eq:thermal-rom}
\end{equation}
这里 \(\bm M\succ0\)，
\(\bm K\) 对应传导/换热耗散，
学习闭合项 \(\bm f_\theta\) 只能表示未建模热通量
或低阶截断，不得任意制造/消灭能量。
在线时间通过对式 \eqref{eq:thermal-rom} 的积分推进，
不是将时间作为单个黑箱网络输入后截断训练窗。
这给出长时间求解的机制，\textbf{并不自动证明长时间精度}；
仍需在未知功率历程、位姿和热稳态下独立测试。

\section{谐振电路与电磁热双向闭环}\label{sec:circuit}
\subsection{实际端口电流}
给定电源/负载、补偿网络拓扑、工作频率和温度状态，
由同一 \(\bm Z_{\rm EM}\) 组合电路阻抗，
在基波线性稳态近似下解
\[
\bm{\mathcal Z}_{\rm circuit}(\bm Z_{\rm EM},\omega,\mathcal C,\bm u)
\bm I=\bm V.
\]
补偿电容、电源和负载不应计入耦合器内部空间热源。
LLC/CLLC 开关器件、整流器和谐波需要频域多谐波或时域模型；
首版只能对明确指定的基波工作域作保证。

\subsection{闭环热时间步}
在热时间 \(t_n\)，以热状态 \(\bm a_n\) 计算材料
\(\mathcal M(T_n)\)，求同源单位端口电磁状态
\((\bm Z_n,\{\bm F_{p,n}\})=\mathcal E_\theta(\mathcal S_n,\mathcal M_n)\)，
解电路电流 \(\bm I_n=\mathcal C(\bm Z_n,\bm u_n)\)，
再由场叠加得到局部 \(q_{i,n}\) 和热动力变化：
\begin{equation}
\boxed{
\bm a_n\longrightarrow \mathcal M(T_n)
\longrightarrow(\bm Z_n,\bm F_n)
\longrightarrow\bm I_n
\longrightarrow q_{i,n}
\longrightarrow\bm a_{n+1}.
}\label{eq:feedback}
\end{equation}
导体电阻、海水电导率和磁性材料参数随温度变化，
可导致阻抗和谐振工况改变。
在温度强反馈情况下应在同一时间步内做隐式/固定点迭代，
弱耦合可采用误差受控的半隐式积分。
只有走通闭环才是双向电磁热代理。

\subsection{慢热/快电分解}
电磁载波周期通常远小于热常数，
故可采用载波相量和慢热包络联立；
启动、负载突变及高频谐波超出包络假设时要切换适当电路瞬态模型。
海水流动和位姿历程也应由真实时间输入，不可在网络中仅附加静态位置标签。

\section{训练任务、物理基准与逐阶段验证}\label{sec:training}
\subsection{四阶段研发}
\begin{enumerate}
\item \textbf{BFZI 解析基线：}直接复核原几何、圆形极限、
跨匝节点仿射性、完整馈回流、两线圈六自由度、
空气电感及误差受控主螺旋源项积分。
\item \textbf{海水复场与分布损耗：}先实现非导电封装、
海水电流连续约束和自洽反作用；从同一单位端口状态
得到 \(\bm Z,\bm E,\bm B,\bm J,q_s,q_{\rm Cu}\)，
再训练海水场散射残差；明确不同材料分区的功率归属。
\item \textbf{连续热场及演化：}构造随体几何条件化基、
界面热通量/对流/海水尾迹和低阶 ODE，
验证功率阶跃、稳态、关断、相对姿态及长时间演化。
\item \textbf{谐振电路双向耦合：}接入确定的电源、
补偿拓扑和负载，比较固定电流/固定电压、
单向/双向热反馈、基波/多谐波模型。
\end{enumerate}
\subsection{学习目标与失效判据}
候选损失
\[
\mathcal L=\lambda_F\mathcal L_F+\lambda_Z\mathcal L_Z+
\lambda_Q\mathcal L_Q+\lambda_T\mathcal L_T+
\lambda_c\mathcal L_{\rm continuity}+
\lambda_e\mathcal L_{\rm energy}+
\lambda_d\mathcal L_{\rm dynamics}.
\]
各项需归一化并报告物理单位误差；
更应采用结构性的散度/耗散/互易约束，不通过堆积损失权重代替。
训练/测试按完整几何、封装、姿态、海水电导率、
频率、功率和时间历程分组；
同一场景的空间点随机分配不是跨几何泛化验证。

\subsection{初始验证工况和未知数据}
可参考 0UnifiedModelingUWPT 的论文草稿：
200/240 mm 线圈、204/244 mm 盒体、
\(N=10,s=6\) mm、\(\alpha=0.5\)、200 kHz、
海水电导率 5 S/m 等工况。
\textbf{其 FEM/实验表格仍有待填内容，不得将其当成已证实性能。}
至少报告端口复阻抗、空间场及热源、最高温度、
长期稳态、能量收支和全场误差；零值附近采用有尺度的绝对误差。
加速比需说明基准、硬件、预计算、加载和批量查询成本。
建议消融：跨匝聚合/逐匝；同源场/端口黑箱；
无海水反作用/有；固定热基/条件化热基；
单向/双向；静止海水/给定海流。

\section{创新假设、风险与下一轮决定}\label{sec:decisions}
候选方法创新是：（1）BFZI 匝号仿射结构同时用于
积分核及神经场源编码；（2）具有海水电流连续约束的非局部场响应学习，
用同一响应计算端口、空间场、损耗；
（3）随体固体温度与世界海水尾迹的双坐标表示；
（4）同源功率一致的热降阶 ODE 与谐振电路反馈。
POD、物理先验与 NN、参数化热 ROM、LLC 电热代理、
神经算子与双向磁热耦合已有先例
\cite{wang2026,dam2026,zhu2026,xie2026,yuan2024,chen2025,dini2023,barba2026}。
仅将这些名词组合不足以证明创新，必须以明确计算复杂度、
误差、跨几何消融和实验结果建立独立贡献。

\textbf{不可声称}任意材料/任意高速运动无需验证、
无固定世界网格等同完全免求积、
端口准确等同空间热源准确、
单频被动性保证全频因果性、
已有论文的速度直接可迁移为本模型速度。

代码实现前应收敛下列五个具体决定：
完整端口和导体截面约定；近场集肤/邻近的基准方法；
海水和导电盒壁材料界面表示；低维电磁状态/投影设计；
几何条件化热基如何保留耗散与跨界面条件。
新理论保留与实验事实的区别，后续以本目录为主线，
不再以旧超椭圆理论并行指导 BFZI 代码。

\section{首版理论封版与进入代码阶段的强制实现契约}\label{sec:freeze}
\subsection{封版决策：从无限可扩展目标收敛到可完成物理闭环}
前面各节提供了可扩展研究框架；此处明确\textbf{第一代码里程碑}
（以下简称 BFZI-MQS-v1）的\textbf{不可再扩大}的基本模型：
\begin{enumerate}
\item 两个电气端口、两根 BFZI 圆形--圆角方形\textbf{连续平面螺旋}，
各自带完全定义的馈/回流路径，参数
\((N,A_0,\varrho_0,s,\alpha,\sigma_w)\) 和真实单位采用 SI；
\item 两侧封装互不穿透，盒体为绝缘、非磁性；海水为
\(\mu=\mu_0\)、各向同性线性正电导率
\(\sigma_s\) 的均匀静止介质；
\item 两侧可取任意\textbf{合法静态}六自由度位姿；时间变化姿态的
准静态逐帧计算可测试，但首版不包括运动电动势与流体运动；
\item 电磁计算为单一实频率 \(\omega=2\pi f\) 的 MQS，
峰值相量约定 \(e^{\jj\omega t}\)；
\item 初始海水教师采用有限海水外域 \(\Omega_B\setminus\Omega_{\rm pack}\)
和绝缘人工外边界 \(n\cdot J_s=0\)，通过增大外域做收敛检验；
\textbf{不得}把此计算冒充严格无限海水解；
\item 首版将 BFZI 线圈作为已给定\textbf{单位端口丝状源}；
空气自项需由 GMD/有限截面规则处理，铜材内部分布损耗暂不预测。
初版 \(Z_{\rm EM}=j\omega L_{\rm air}+Z_{\rm sw}\)
只含线源磁响应和海水耗散，不能宣称总导体阻抗；
\item 仅当电磁端口、空间海水电流和耗散均来自同一可重现
REFERENCE 状态，才开始低阶神经子空间训练；
热与电路模块用同一个场损耗投影作为后续里程碑，暂不阻挡
BFZI-MQS-v1 开工。
\end{enumerate}
上述限制是\textbf{首版实现边界}而不是最终系统能力边界。
本轮起理论范围冻结，今后只允许因实现发现的数学错误修正，
不得为等待所有高级介质和热流细节而延迟主代码。

\subsection{实际闭合的 REFERENCE 算法}
对每个端口 \(p\) 用 BFZI 解析直边/圆角路径和
馈/回流段建立 \(\bm C_p\)；对空气感应项用
Neumann 分组求积、有限截面自项约定建立
\(\bm L_{\rm air}\)。对有限海水域先使用具有
\(H(\mathrm{div})\) 连续性、绝缘界面零法向通量的
\textbf{确定性局部 curl/Piola 字典} \(\bm W_0(\bm x;\mathcal G)\)，
明确包括远场及当前拓扑所需模式。
\(\bm R,\bm L,\bm K\) 按
\eqref{eq:Rs}--\eqref{eq:Kbfzi} 的原始体积分/Green 双体积分直接组装；
先不强制 Fourier 特征加速和网络。
\[
\bm A_s=\bm R+\jj\omega\bm L,\qquad
\bm X_p=-\jj\omega\bm A_s^{-1}\bm K\bm e_p.
\]
在 \(\bm R\succ0\) 时
\[
\bm Z_{\rm sw}=\omega^2\bm K^{\mathsf T}\bm A_s^{-1}\bm K,
\quad
q_s(\bm x;\bm I)
=\frac{\|\bm W_0(\bm x)\bm X\bm I\|^2}{2\sigma_s},
\quad
\bm Z_{\rm EM}=\jj\omega\bm L_{\rm air}+\bm Z_{\rm sw}.
\]
输出单位端口基状态 \(\bm X\in\mathbb C^{M\times2}\)，
此后可以对任何 \(\bm I\) 组合场/耗散而不重复建立基。
\textbf{若局部字典在某几何下覆盖不足，应增广字典并对比参考教师；
内部代数功率一致性不等于真实海水涡流场已收敛。}

\subsection{首轮 FAST 算法及其严格解码}
首版 FAST 不采用独立端口 NN。先运行\textbf{确定性}
低秩基选择/投影，建立无网络的子空间闭环：
\[
\bm W_r=\bm W_0\bm S,\quad
(\bm R_r,\bm L_r,\bm K_r)
=(\bm S^{\mathsf T}\bm R\bm S,
\bm S^{\mathsf T}\bm L\bm S,\bm S^{\mathsf T}\bm K).
\]
先由 SVD/贪心或能量加权正交方法从教师单位端口响应
选择 \(\bm S\)；在线低阶矩阵必须仍按同一物理积分定义。
通过验证后再用轻量条件网络替代 \(\bm S\) 的选择策略
（而不是让 NN 自由生成任意 \(Z,q\)）。
\textbf{可测加速的前提}：矩阵/核的参数化复用和 BFZI
匝号聚合的总耗时少于参考组装；
若每次仍建立全 \(M\times M\) 矩阵，不能称为高性能 FAST。
此项性能工程在代码阶段对直接求积、特征预计算和
BFZI 块积分逐项计时后选择，不再阻碍物理 REFERENCE 首版。

\subsection{物理接口的数据和相量合同}
场景参数不得把 BFZI 的 \(\varrho_0\) 误当成超椭圆形状阶数。
端口用指定\textbf{峰值}电流方向，并定义完整馈线/回流路径，
不允许仅依靠无限开口曲线声称完整实际端口自感。
固定输出：
\[
\begin{aligned}
\mathrm{EMPortResponse}&=\{f,\bm Z_{\rm EM},
\bm Z_{\rm sw},\bm L_{\rm air},
\mathrm{validity},\mathrm{reference\_metadata}\},\\
\mathrm{EMFieldResponse}&=\{\bm x,\mathrm{region},
\bm E(\bm x),\bm J_s(\bm x),\bm B(\bm x),q_s(\bm x)\},\\
\mathrm{ROMState}&=\{\bm W_r,\bm R_r,\bm L_r,\bm K_r,
\bm X_r,\mathrm{scene\_key}\}.
\end{aligned}
\]
在静态线性条件下，\(\bm E,\bm J,\bm B\) 输出可表示为
两个单位端口基场的线性组合；
\(q_s\) 按 \(|\bm J_s|^2/(2\sigma_s)\) 非负二次计算，
不能把单位端口热源简单线性相加
（两个端口之间有交叉项）。
端口实部损耗归属与 \(q_s\) 的积分采用\textbf{一致海水状态}；
在有限积分误差下报告数值闭合残差。
若同时请求不同几何、温度或频率，
必须重建或更新相应物理矩阵，不可错用缓存。

\subsection{实现顺序与代码包边界}
考虑到原 \texttt{sdfmpneo\_vnext} 实现体量较大
且已有 \texttt{SuperellipseSpiral} 代表另一种形状族，
新实现应在\textbf{同一仓库的独立 Python 子包}
\texttt{sdfmpneo\_bfzi/} 中起步，再经适配器接入
\texttt{MeshfreeVNextSystem} 和训练入口。
安装配置中 \texttt{include=["sdfmpneo*"]} 能发现该子包。
第一批按依赖方向实现：
\begin{enumerate}
\item \texttt{geometry.py}：BFZI 参数验证、
八线段组 \((\bm r,\bm g)\)、连续圆形退化、
刚体位姿与 Gauss 节点、馈/回流几何接口；
\item \texttt{scene.py}：两个不交绝缘封装、海水导电率、
人工外域、区域分类和合规几何检查；
\item \texttt{source.py}：BFZI 完整端口的
\(\bm C_p\) 与空气电感积分，先逐匝参考再接
有误差控制的匝号快速块；
\item \texttt{sea\_basis.py}：确定性 curl 局部图谱、
零法向模式、远场/间隙字典，按真实场景输出评估器；
\item \texttt{sea\_reference.py}：正权 \(\bm R\)、
一致 \(\bm L\) 和 \(\bm K\) 参考组装、
两端口单位激励求解、场读出与逐材料海水热源；
\item \texttt{sea\_rom.py}：确定性低阶投影，
保结构 \(\bm Z_{\rm sw}\) 与非负空间热源；
\item 后续 \texttt{thermal.py},\texttt{circuit.py},\texttt{neural.py}
接入空间损耗、双向电热演化和 NN 子空间。
\end{enumerate}
第一阶段\textbf{不迁移历史训练脚本、不实现新 GUI、
不改旧 \texttt{SuperellipseSpiral} 定义、不建立 GitHub Actions}。

\subsection{完成 BFZI-MQS-v1 的唯一验收门槛}
代码阶段针对少量自洽物理算例直接做以下必要检查：
\begin{enumerate}
\item BFZI 主曲线的直边--圆角段连接、
有向线元、圆形极限、任意刚体位姿下的仿射匝号保持，
以及有限截面自项约定；
\item 对双盒合法姿态，海水模式在人工外边界与绝缘封装面
满足零法向电流，且求积电阻矩阵满秩；
\item 同一海水离散空间内
\(\bm Z_{\rm sw}\) 数值对称、Hermitian 部分非负，
\(\int q_s\,\mathrm dV\) 与
\(\tfrac12\bm I^{\mathrm H}\operatorname{Herm}(\bm Z_{\rm sw})\bm I\)
在相同积分规则下闭合；
\item 随外域增大、求积与电流基增广，
端口结果及局部场向稳定参考收敛；
\textbf{只有离散功率恒等式而没有此收敛不能称作参考精度达标}；
\item 空气/海水基准与原 BFZI 公式一致，
且完整推理耗时记录包括节点生成、矩阵组装与求解。
\end{enumerate}
这些是主体功能的最小验收，不需要衍生大量审计/认证框架。
\textbf{进入代码的决定：已满足。}
后续复杂材料、全域神经模型、热动力及电路闭环作为代码迭代工作，
不再提出额外的理论封版前置条件。

% G1 theory-only correction and verification contract (2026-11-04 UTC).
% Restricts claims to the linear MQS baseline. Not an assertion that G1 has passed.
\section{G1 基础方程、近似空间与可信参考解的重新核查}
\label{sec:g1-foundation-recheck}

\subsection{首先区分三个不同的数学对象}
G1 最终要近似的是无限均匀导电海水中的连续磁准静态问题。
设两个互不相交的绝缘实心封装为 \(B_1,B_2\)，
\[
\Omega_\infty=\RR^3\setminus(\overline{B_1}\cup\overline{B_2}),
\qquad
\Omega_R=B_R(\bm c_0)\setminus(\overline{B_1}\cup\overline{B_2}).
\]
第一层是 \(\Omega_\infty\) 上满足无穷远物理衰减的场；
第二层是 \(\Omega_R\) 上指定了\emph{人工}外边界条件的连续问题；
第三层才是有限维试探函数及数值积分得到的离散问题。
\textbf{有限维功率闭合、有限外域相邻解接近、磁核矩阵半正定，
分别不能证明这三层问题已经一致。}

首版只假设静止、线性、非磁性绝缘封装、
\(\mu=\mu_0\)、均匀正海水电导率 \(\sigma_s\)、
闭合的外加 BFZI 两端口电流源及峰值相量 \(e^{\jj\omega t}\)。
不把本模型扩大解释为含磁芯非线性、导电屏蔽、
位移电流、运动感应和铜线内全耦合涡流的通用参考解。

\subsection{从 Maxwell 方程严格得到电流弱式}
在海水无外加体电流的区域，磁准静态方程为
\begin{align}
\Curl\bm H&=\bm J_s,&
\Curl\bm E&=-\jj\omega\bm B,&
\Div\bm B&=0,\\
\bm B&=\mu_0\bm H,&
\bm J_s&=\sigma_s\bm E,&
\Div\bm J_s&=0.
\end{align}
完整闭合源线电流的单位磁矢势记为
\(\bm A_{{\rm src},p}=\mu_0\bm C_p\)，
\(\bm C_p=\frac1{4\pi}\int_{\Gamma_p}\mathrm d\bm\ell'/|\bm x-\bm r'|\)。
取合适规范下的自由空间核
\((\mathcal G_0\bm J)(\bm x)=\int_{\Omega}\bm J(\bm y)/
(4\pi|\bm x-\bm y|)\,\mathrm dV_y\)，则
\[
\sigma_s^{-1}\bm J_s+\jj\omega\mu_0\mathcal G_0\bm J_s
+\grad\phi=-\jj\omega\mu_0\sum_p\bm C_p I_p.
\]
对有限外域，若使用人工绝缘截断条件，正确的测试空间为
\begin{equation}
\boxed{
\mathcal V_R=\{\bm v\in H(\operatorname{div};\Omega_R):
\Div\bm v=0,\quad
\bm v\cdot\bm n=0\ \text{于 }
\partial B_1\cup\partial B_2\cup\partial B_R\}.
}
\label{eq:g1-full-boundary-space}
\end{equation}
关键是\textbf{整个有限域边界}都要出现在零法向条件中；
否则
\[
\int_{\Omega_R}\overline{\bm v}\cdot\grad\phi\,\mathrm dV
=\int_{\partial\Omega_R}\phi\,\overline{\bm v}\cdot\bm n\,\mathrm dS
\]
在人工外球处未必为零，不能从弱式中凭空删除该边界项。
如采用非反射人工边界、允许非零外球电流通量，则必须
\textbf{保留并正确处理}对应边界算子/边界功项，
不能套用同一个无边界项弱式。
无限域则应使用具有相应远场衰减和磁能可积性的函数空间，
证明无穷远边界项消失后再作消元。

对所有 \(\bm v\in\mathcal V_R\) 得
\begin{equation}
a_R(\bm J_s,\bm v)
=\int_{\Omega_R}\sigma_s^{-1}\bm J_s\cdot\overline{\bm v}\,\mathrm dV
+\jj\omega\mu_0\int_{\Omega_R}
(\mathcal G_0\bm J_s)\cdot\overline{\bm v}\,\mathrm dV
=-\jj\omega\sum_p I_p\,(\bm A_{{\rm src},p},\bm v)_{L^2}.
\label{eq:g1-weak-corrected}
\end{equation}
这是\emph{人工绝缘有限域问题}的正确变分形式，
并不表示该问题等于无限海水的限制。

\subsection{哪些有限维结构结论确实正确}
选取实值线性独立的 \(\bm W_m\in\mathcal V_R\)，
以本稿式 \eqref{eq:Rs}--\eqref{eq:Kbfzi}
精确或一致求积定义 \(\bm R,\bm L,\bm K\)。
对实正 \(\sigma_s\)、自由空间对称 Coulomb 核及闭合端口，
理论上
\[
\bm R=\bm R^{\mathsf T}\succ0,\quad
\bm L=\bm L^{\mathsf T}\succeq0,\quad
(\bm R+\jj\omega\bm L)\bm z=-\jj\omega\bm K\bm I.
\]
令 \(\bm A=\bm R+\jj\omega\bm L\)，则
\[
\bm Z_{\rm sw}
=\omega^2\bm K^{\mathsf T}\bm A^{-1}\bm K,\qquad
\operatorname{Herm}(\bm Z_{\rm sw})
=\omega^2\bm K^{\mathsf T}\bm A^{-\mathrm H}
\bm R\bm A^{-1}\bm K\succeq0.
\]
故 \(\bm Z_{\rm sw}^{\mathsf T}=\bm Z_{\rm sw}\)，且
\[
\tfrac12\bm I^{\mathrm H}\operatorname{Herm}(\bm Z_{\rm sw})\bm I
=\tfrac12\bm z^{\mathrm H}\bm R\bm z
=\int_{\Omega_R}\frac{|\bm J_s|^2}{2\sigma_s}\,\mathrm dV.
\]
这些是\textbf{同一个离散系统的代数定理}；
它们证明约束、互易和耗散结构，不给出
\(\|\bm J_s-\bm J_\infty\|\) 或热点误差的任何界。
以离散自身的 \(\bm R,\bm K\) 重建功率，
即使 \(\bm L\) 或电流空间错误，仍可在机器精度闭合。

\subsection{当前 curl--高斯基的数学优势与缺口}
实现取 \(\bm W_m=\Curl\bm F_m\)，
\(\bm F_m\) 是高斯包络、绝缘盒体距离因子及外球截断因子
的乘积再乘常向量或完整端口的源磁矢势。
在适当的分片 \(H(\operatorname{curl})\) 正则性及相容迹下，
若 \(\bm n\times\bm F_m=0\)，可保证
\(\Div\bm W_m=0\) 与 \(\bm n\cdot\bm W_m=0\)；
这是本方案可保留的正确设计。

但当前外球 \(C^2\) 截断函数同时满足
\(f(R)=f'(R)=0\)，因此在外球处实际强制
\(\bm W_m|_{\partial B_R}=0\)，不仅是法向电流为零。
这在有限模态下是\emph{过于受限的合规子空间}，
并不等价于整个 \(\mathcal V_R\)；
它在 \(L^2\)/\(H(\operatorname{div})\) 完备性层面
未必构成不可逾越的阻碍，但需要\textbf{系统的空间加密证明}，
不能由 229 个高斯方向或有限几种长度尺度自行推出。
改变外球半径还会改变所有模式的边界包络，
故即使每个高斯中心、宽度及方向保持不变，
\textbf{跨半径函数空间也不是严格嵌套的}；
相邻解差异同时包含外域、外球边界和试探空间改变。

两个不相交、各自同胚于三维闭球的\emph{实心}绝缘盒体
在三维球域中的补集，一般没有一维同调环的拓扑障碍；
\textbf{不能由“有两个封装/多个边界分支”直接推断
必须加入非零的一阶拓扑谐和电流}。
物理上仍需表示跨封装及海水远场的\emph{普通大尺度涡流}，
这是空间逼近问题而不是自动出现的拓扑缺陷。
几何若改成穿孔实体、导电回路等，才需要重新计算拓扑条件。

\subsection{磁 Green 离散与连续 Coulomb 双积分不可混淆}
目标磁算子是
\[
L_{mn}
=\frac{\mu_0}{4\pi}
\int_{\Omega_R}\int_{\Omega_R}
\frac{\bm W_m(\bm x)\cdot\bm W_n(\bm y)}
{|\bm x-\bm y|}\,\mathrm dV_y\mathrm dV_x.
\]
目前 \texttt{uniform\_ball} 规则把每个求积权重视作等体积
\emph{球形电流云}，用球--球平均 Coulomb 核替代真实单元内的
双体积分。这能构造对称非负的近似磁能，
但球未必是实际正权积分单元的形状或支撑，
也未必保持真实海水边界和点值电流的无散度结构；
\textbf{正定性不蕴含对上述目标算子的数值一致性}。
必须独立检查采样节点分辨率、球形替代误差、
近对角奇异项与非局部尾项；当模式数和外域改变时不能
仅复用低精度 Green 矩阵并宣称参考级精度。
当前 229 模态 A/B \(\|\Delta L\|_F/\|L_A\|_F\approx13.2\%\)，
确实没有算子收敛证据。
这一 Frobenius 相对量本身也不能代替
\(\omega\|\bm R^{-1/2}\Delta\bm L\bm R^{-1/2}\|_2\)
或被激发方向的残差响应敏感度；
需同时检查这些\emph{无量纲动力学相关指标}，
不能仅由一个全矩阵百分比判断各端口场的误差。

\subsection{实际海水穿透深度与截断半径}
对于远离封装的均匀导电海水，在 MQS 假设下
\[
\grad^2\bm H=\jj\omega\mu_0\sigma_s\bm H,\qquad
\gamma=\sqrt{\jj\omega\mu_0\sigma_s}
=\frac{1+\jj}{\delta_s},\qquad
\delta_s=\sqrt{\frac{2}{\omega\mu_0\sigma_s}}.
\]
默认示例 \(f=200\,\mathrm{kHz}\)，
\(\sigma_s=5\,\mathrm{S/m}\)，
\(\mu=\mu_0\) 时 \(\delta_s\simeq0.503\,\mathrm m\)。
目前 \(0.52\)--\(1.30\,\mathrm m\) 的球半径仅约
\(1.03\)--\(2.58\) 个穿透深度（从球心计）；
因此外域远场电流及其磁反馈未必足够小，
\textbf{不能预设这些半径应形成收敛平台}。
球形/非均匀几何中的振幅还含角向及代数衰减，
实际截断误差不能只用 \(e^{-R/\delta_s}\) 代替测量。
较大的球域也会增加 Green 积分的求积误差：
外域相邻解非单调并不能单凭日志判断是物理尾部还是积分别名。

\subsection{从代数残差升级到 G1 可解释的误差界}
在固定\emph{连续}有限域问题中定义
\(\|\bm u\|_{\mathcal R}^2=
\int_{\Omega_R}\sigma_s^{-1}|\bm u|^2\,\mathrm dV\)。
因 \(\operatorname{Re}a_R(\bm u,\bm u)
=\|\bm u\|_{\mathcal R}^2\)，
对真实有限域解 \(\bm J_R\) 和任意合规近似
\(\bm J_h\)，其在\emph{完整测试空间}上的残差
\[
\mathfrak r_h(\bm v)
=\ell_R(\bm v)-a_R(\bm J_h,\bm v),\qquad
\|\mathfrak r_h\|_{\mathcal V_R'}
=\sup_{\bm v\in\mathcal V_R\setminus\{0\}}
\frac{|\mathfrak r_h(\bm v)|}{\|\bm v\|_{\mathcal R}}
\]
满足
\begin{equation}
\boxed{\|\bm J_R-\bm J_h\|_{\mathcal R}
\leq\|\mathfrak r_h\|_{\mathcal V_R'}.}
\label{eq:g1-residual-continuum}
\end{equation}
该结论是对\emph{相同连续算子与完整空间}的有条件估计，
不是当前离散 229 父空间上的残差上界。
只在另一个有限测试字典上取 supremum 时，
得到的是完整对偶残差范数的\emph{下界}，
除非再建立可靠的稳定性、饱和性或误差上界估计。
因此目前 \texttt{parent\_dJ\_energy\_upper}
只证明\emph{给定父空间内部}的逼近效果，
不能作无限维电磁参考误差认证。

对空间焦耳热有连续积分控制
\[
\|q(\bm J)-q(\bm J_h)\|_{L^1(\Omega_R)}
\leq\tfrac12\|\bm J-\bm J_h\|_{\mathcal R}
\big(\|\bm J\|_{\mathcal R}+\|\bm J_h\|_{\mathcal R}\big),
\]
但这并不控制盒壁极局部的 \(L^\infty\) 热点；
必须另外检查几何随体的壁面/边缘/角部点列及更强局部范数。

\subsection{第一批可实施的数学重构（小空间实验，而非直接认证）}
\paragraph{不修改 Coulomb 核的近场奇异分拆}
为避免将任意高斯求积权重伪装成\emph{球形体积单元}，
固定海水测试空间和有限域，取
\(\bm x\in\Omega_R\)，
\(0<a(\bm x)<\operatorname{dist}(\bm x,\partial\Omega_R)\)。
构造 \(C^2\) 径向窗 \(\chi_a(\rho)\)，
在 \(\rho\le a/2\) 时为 \(1\)，在 \(\rho\ge a\) 时为 \(0\)。
则对任何可积向量场 \(\bm W_n\) 有严格恒等式
\begin{align}
\int_{\Omega_R}\frac{\bm W_n(\bm y)}{|\bm x-\bm y|}\,\mathrm dV_y
={}&\int_{\Omega_R}
\frac{1-\chi_a(|\bm x-\bm y|)}
{|\bm x-\bm y|}\bm W_n(\bm y)\,\mathrm dV_y
\nonumber\\
&+\int_{S^2}\int_0^{a(\bm x)}
\rho\,\chi_a(\rho)\,
\bm W_n(\bm x+\rho\bm s)\,\mathrm d\rho\,\mathrm d\Omega_{\bm s}.
\label{eq:g1-coulomb-exact-split}
\end{align}
第二项在物理球坐标中解析消除了 \(1/\rho\) 奇异因子，
\emph{不引入核软化或节点云球模型}。
这是核恒等式，离散计算仍须分别对
\(\bm x\)、远场 \(\bm y\)、近场径向和角度作独立加密，
并保证相应求积的积分区域与模式评估精确一致。
实现 \texttt{sdfmpneo\_bfzi/green\_coulomb.py} 是有限模态小型
对照器，不能直接宣称其粗节点解精确或可快速运行 229 模态；
使用不同 \(\bm x/\bm y\) 求积时暂以对称化矩阵诊断，
\textbf{不能强迫其正定以掩盖积分误差}。

\paragraph{可逐级加密的紧支 curl 图谱}
用 \(C^2\) Wendland 核
\[
\psi(s)=(1-s)_+^4(4s+1),\qquad
\grad\psi(\sqrt{\bm r^{\mathsf T}G\bm r})
=-20(1-s)_+^3G\bm r
\]
替换局部高斯分布的固定有限个采样方向，
在封装外的面片以 \(2^\ell\times 2^\ell\)
网格状\emph{中心选点}生成紧支向量势，
再取 \(\Curl(\psi\,d_{\rm ins}\,\bm a)\)。
这不是有限元体网格，而是\emph{解析紧支基函数的层级索引}，
可在层级加密时保持所有已存在的电流方向与
精确解析 curl。额外建立按物理封装长度而非人工球半径
选择中心的远场图谱。实现
\texttt{build\_sea\_basis(structured\_level=1,2)} 已构造此嵌套字典；
目前尚无空间稠密性定理或跨几何误差界。

\paragraph{真正无限域的解析边界基准}
把任意 BFZI 双盒问题直接附加一个经验 Robin 常数是错误的。
先在最简单、可解析的绝缘球腔内中心磁偶极模型中求解。
对偶极的轴对称环向模 \(A_\phi=f(r)\sin\theta\)，
令 \(\gamma^2=\jj\omega\mu_0\sigma_s\)、
\(\operatorname{Re}\gamma>0\)，则海水中
\[
f''+\frac2r f'-\frac2{r^2}f-\gamma^2 f=0,\qquad
f_\infty(r)=Ck_1(\gamma r),
\]
\(k_1\) 为衰减的修正球 Bessel 函数。
在任何包围球腔的半径 \(R\)，此\emph{唯一球谐扇区}
的精确无限域透明条件为
\begin{equation}
\boxed{\frac{f_\infty'(R)}{f_\infty(R)}
=\gamma\frac{k_1'(\gamma R)}{k_1(\gamma R)}.}
\label{eq:g1-exact-dtn-l1}
\end{equation}
内侧真空向量势由 \(M/r^2+Dr\) 组成；
在腔壁连续匹配 \(f\) 和 \(f'\)，
可以得到不含人工外球的精确海水 \(J_\phi=-\jj\omega\sigma_s
f_\infty\sin\theta\) 与 \(q=|J_\phi|^2/(2\sigma_s)\)。
这已在 \texttt{sdfmpneo\_bfzi/exterior\_dipole.py} 中实现，
同时可比较把外球的 \(A_\phi\) 强制为零导致的截断误差。
\textbf{式 \eqref{eq:g1-exact-dtn-l1} 不适用于任意 BFZI
线圈、多极混合和两个长方体封装。}
推广须处理完整矢量球谐的极向/环向模及耦合边界映射，
或改用严谨的边界积分施加非局部透明条件。

\subsection{真实封装几何裁剪的固定半径 Coulomb 奇异补偿}
第一版为保证近场球全部落在海水，使用
\(a(\bm x)=\min(a_0,0.6\,\mathrm{dist}(\bm x,\partial\Omega_R))\)。
在封装面附近，这会使 \(a(\bm x)\to0\)；
理论恒等式并未失效，但需要由粗全局 \(\bm y\) 求积
解析越来越尖锐的近对角核，易出现积分次序别名与相消。
现在允许固定物理补偿半径 \(a_0>0\)，
但\emph{严格按真实海水域裁剪}球形补偿区：
\[
\int_{\Omega_R\cap B_{a_0}(\bm x)}
\frac{\chi_{a_0}(|\bm y-\bm x|)\bm W_n(\bm y)}
{|\bm y-\bm x|}\,\mathrm dV_y
=\int_{S^2}\sum_{[s_k,t_k]\subset I_{\bm x,\bm n}}
\int_{s_k}^{t_k}\rho\,\chi_{a_0}(\rho)
\bm W_n(\bm x+\rho\bm n)\,\mathrm d\rho\,\mathrm d\Omega,
\]
其中 \(I_{\bm x,\bm n}\) 是沿固定射线
\(\bm x+\rho\bm n\) 的海水区间与 \([0,a_0]\) 的交，
用真实旋转盒的射线--板交点和外球交点构造，
不在封装内部评估任何电流模态。
全局 \(\bm y\) 部分仍使用
\((1-\chi_{a_0})/|\bm x-\bm y|\)，
因此\emph{连续}双积分依旧严格等于原始 Coulomb 算子。

此修正仅移除了\emph{近壁补偿球被迫收缩}的结构性限制，
并未使数值求积自动收敛。因为球被盒体切割，
沿方向 \(\bm n\) 的积分具有遮挡/切触轮廓，
角向采样也必须独立加密；必须同时检查
\(a_0\)、局部角向阶数、\(\bm x/\bm y\) 体积规则，
以及对驱动相关的 \(\omega R^{-1/2}LR^{-1/2}\) 的影响。
2026-11-04 的旧规则实测
\(dL_{\bm x}=9.8049\%\)、\(dL_{\bm y}=3.4771\%\)、
\(dL_{\bm x\bm y}=2.9804\%\)，
非加性指标 \(10.4896\%\)，
不能仅据双轴误差抵消宣称新旧任何一方已收敛。

\paragraph{针对近壁方向非光滑性的局部求积修正（新代码待实测）}
2026-11-04 的真实几何裁剪试验已经证明，
近壁补偿球不能只依赖均匀角向规则：
近区角向 \(4/8\to6/12\) 的 \(dL\) 从旧缩球的
\(4.51\times 10^{-7}\) 增大到 \(3.60734\times10^{-3}\)，
改变补偿半径 \(a:0.055\to0.080\,\mathrm m\)
使 \(dL=2.11519\%\)。
独立加密双重积分的外层 \(\bm x\) 给出
\(dL_{\bm x}=11.2368\%\)，内层 \(\bm y\) 仅
\(1.37432\%\)；两轴同时变化 \(1.14953\%\)
来自明显的非加性，并非物理精度通过。

改进方法不改变原 Coulomb 核、实际海水区域、
基函数或局部补偿球：
对最邻近的一处盒壁/人工球面选择局部球面极轴
\(\bm n\)，并将角向积分划分为两段 Gauss 面板。
距无限平面盒壁正距离 \(d<a\) 时，第一处半径
\(a\) 射线触壁对应 \(\mu_*=\cos\theta_*=d/a\)，
取面向固体的法向为极轴。
对于中心 \(\bm c_0\)、半径 \(R\) 的人工外球，
\[
\mu_*=\frac{R^2-|\bm x-\bm c_0|^2-a^2}
{2a|\bm x-\bm c_0|},
\]
极轴沿球心指向 \(\bm x\)。
两个面板 \([-1,\mu_*]\) 和 \([\mu_*,1]\)
各自使用独立的 Gauss--Legendre 极角积分，
方位角使用周期求积；
截断盒角/盒边以及多盒遮挡仍由射线交点计算，
此角向坐标选择\emph{不改变}真实湿区积分。
与此同时，在径向窗 \(\chi_a\) 的分片连接点
\(\rho=a/2\) 显式切开每一个湿区区间，
减少单个低阶 Gauss 规则穿越 \(C^2\) 分片连接
时的高阶导数非光滑性。

平面绝缘壁常向量场有可复核的独立基准：
令采样点距壁 \(d\)，局部球半径 \(a>d\)，则
\[
\int_{\Omega_{\rm half-space}\cap B_a(\bm x)}
\frac{\chi_a(|\bm y-\bm x|)}
{|\bm y-\bm x|}\,\mathrm dV_y
=\int_0^d 4\pi\rho\chi_a(\rho)\,\mathrm d\rho
+\int_d^a 2\pi(\rho+d)\chi_a(\rho)\,\mathrm d\rho .
\]
此积分在固定角向规则下可以直接比较旧的单极角
与新的面向实体的分段极角，且不是人工制造的
原核软化或拟合标签。
它\emph{仅}验证局部半空间的角向裁剪；
G1 仍要求真实旋转双盒场景的
\(\bm x/\bm y\) 求积、分割半径、近壁局部热及
驱动方向稳定性，并检查不同 \(\ell_{\max}\)
的远区物理匹配。

\subsection{均匀导电海水的完整球谐外区映射：当前实现边界}
选择包围所有绝缘封装、完整端口源的球面
\(S_R\)，球外只有均匀正标量电导率海水。
由 \(\Div\bm J=0\)、
\(\Curl\bm E=-\jj\omega\mu_0\bm H\)、
\(\Curl\bm H=\bm J=\sigma_s\bm E\) 得
\[
(\Delta-\gamma^2)\bm J=0,\qquad
\gamma^2=\jj\omega\mu_0\sigma_s,\quad\Re\gamma>0.
\]
故在固定 Cartesian 分量下，对\emph{散度相容}的边界迹
\(\bm J(R,\theta,\varphi)\)，外区衰减解为
\begin{equation}
J_i(r,\theta,\varphi)
=\sum_{\ell=0}^{\infty}\sum_{m=-\ell}^{\ell}
j_{i,\ell m}(R)
\frac{k_\ell(\gamma r)}{k_\ell(\gamma R)}
Y_{\ell m}(\theta,\varphi),\qquad r\ge R.
\end{equation}
球面透明算子是角向非局部的
\begin{equation}
\partial_r J_i(R,\theta,\varphi)
=\sum_{\ell,m}
\underbrace{\gamma\frac{k'_\ell(\gamma R)}
{k_\ell(\gamma R)}}_{\lambda_\ell(R)}
j_{i,\ell m}(R)Y_{\ell m}(\theta,\varphi).
\label{eq:g1-multipolar-dtn}
\end{equation}
这解释了为什么\textbf{只有一个} Robin 系数在一般
双盒 BFZI 问题中不正确，且其 \(\ell=1\)
子块与前述解析偶极基准相符。
利用球谐正交性，已知相容球面电流系数时，外海水
平均 Joule 功率为
\[
P_{\mathrm{tail}}
=\frac1{2\sigma_s}\sum_{i,\ell,m}
|j_{i,\ell m}(R)|^2
\int_R^\infty r^2\left|
\frac{k_\ell(\gamma r)}{k_\ell(\gamma R)}
\right|^2\,\mathrm dr .
\]

实现 \texttt{ExteriorSphericalCurrent} 提供
有限 \(\ell_{\max}\) 截断下的投影、外区续接、
复对数导数与尾区损耗，
并加入 \(\ell=1,2\) 环向无散度模态数值基准。
\textbf{关键未完成：}
从有限域原始 Coulomb Galerkin 解得到的任意球面迹
未必满足无限域散度相容条件，且现有外球包络将
边界电流强制压零。必须将式
\eqref{eq:g1-multipolar-dtn}
作为非局部边界条件\emph{反向耦合}到双盒内区求解，
并控制角向截断、匹配球半径与真物理边界误差。
只对某个已截断的数值场做球谐外推，
不构成完整 BFZI 无穷域误差证书。

\subsection{纠正后的 G1 通过条件与工作次序}
\begin{enumerate}
\item \textbf{正确边值问题：}
证明每个频率与材料参数下源回路闭合，
说明无限域衰减与人工外域边界算子的关系，
区分非导电封装的零法向和人为外球条件；
检查 MQS 适用范围（例如 \(\omega\epsilon/\sigma\ll1\)）。
\item \textbf{真实磁核：}
在固定物理基及同一个外域上，独立加密 Green 积分、
控制对角近奇异项、球云形状替代和远场尾项，
给出动力学加权磁算子及双端口受迫场的一致误差，
而非仅验证 PSD。
\item \textbf{有系统的海水空间逼近：}
建立随尺度加密的紧支 curl/Piola 近盒、间隙、
完整远场角向/径向字典，保证合规与适当稠密性；
对独立大父空间评估 J、双端口 J、q 和残差，
不是只在固定 229 父解里做压缩。
\item \textbf{外域独立收敛：}
固定共同物理湿区，同时用真正物理一致的
外域尾场/透明边界或多个充分大的截断半径核对；
必须把外域截断、边界包络变化、R/K 求积和 Green L
积分误差作为\emph{不同}来源逐项控制。
\item \textbf{独立物理基准：}
至少在代表性几何、错位姿态和封装尺寸下与
不同离散原则的高精度解（解析可用极限、
独立边界积分或其他可信参考）交叉核对，
报告端口、海水体积 J、封装近壁 q 和功率误差。
允许参考方法使用局部离散；
主代理仍可坚持无固定世界网格设计。
\end{enumerate}
在以上各项证据成立以前，
\textbf{G1=未通过，229 模态不能用作正式跨几何神经教师}。
下一轮优先分离同一物理坐标中的 R/K 与 Green L
敏感性，接着修磁核求积和边界/空间逼近；
不再把简单扩大外域或增加少数经验高斯模式
当作充分的 G1 收敛方案。

% Independent theory/numerics review, 2026-11-04.
% The section is a contract for future verification, NOT a convergence claim.
\section{G1 理论与数值积分的文献对照及重新验收顺序}
\label{sec:g1-numerics-literature}

\subsection{确认 Maxwell--MQS 问题而非仅确认 Galerkin 矩阵}
取两只不相交的密封绝缘实心盒 $B_1,B_2$、
$\Omega_\infty=\RR^3\setminus(\overline B_1\cup\overline B_2)$，
静止、线性、均匀 $\mu=\mu_0$、$\sigma_s>0$，
闭合 BFZI 端口路径及相量 $e^{\jj\omega t}$。
必须同时检查 MQS 所忽略的位移电流、传播长度、
线源的散度兼容性与导电域拓扑条件。
Ammari--Buffa--N\'ed\'elec 对涡流近似的源兼容
及低频误差阶有明确分析~\cite{ammaribuffanedelec2000}。
仅给出正定的 $\bm R$ 和被动的阻抗，
\emph{不能反向证明} Maxwell 模型和外域处理正确。

将符合绝缘封装零法向和无穷远衰减的海水涡流
记为 $\bm J$。在合适连续空间内，
\[
\bm A_{\rm src}=\mu_0\sum_p\bm C_pI_p,\quad
\bm A_{\rm sea}=\mu_0\mathcal G_0\bm J,\quad
\sigma_s^{-1}\bm J+\jj\omega(\bm A_{\rm src}+
\bm A_{\rm sea})+\grad\phi=0,
\]
其中 $\mathcal G_0\bm J(\bm x)=\int_{\Omega}
\bm J(\bm y)/(4\pi|\bm x-\bm y|)\,\mathrm d\bm y$。
对所有散度为零且全边界法向迹为零的测试场 $\bm v$，
标量势项才能通过分部积分消去。
由此所得的
$\bm R+\jj\omega\bm L$、$\bm K$、$\bm Z_{\rm sw}$
及耗散恒等式已在
第~\ref{sec:galerkin} 和 \ref{sec:g1-foundation-recheck} 节证明；
它们是\textbf{正确的条件代数结构}，
不是原 Coulomb 积分或实际电流近似的误差证书。

\subsection{电流的合法性、完备性和人工外球}
令 $\bm W_m=\Curl\bm F_m$，
若 $\bm F_m$ 有合适的 $H(\mathrm{curl})$ 正则性且
$\bm n\times\bm F_m=0$，
则 $\bm W_m$ 无散度、零法向迹。
但这不能推出有限高斯／紧支字典在目标
$H(\mathrm{div})$ 电流空间中的稠密性；
需建立插值/近似、稳定性和真正的层级收敛性。
符合微分形式的空间构造及收敛问题见
Hiptmair~\cite{hiptmair2002}；
局部环流与可能的拓扑自由度的不同职责
见 Voltolina 等~\cite{voltolina2021}。
两只不穿孔的实心盒本身不能自动推出一阶环流拓扑缺陷。

当前外球的窗函数通常令边界
$\bm W_m|_{\partial B_R}=0$，
比人工绝缘边界的 $\bm n\cdot\bm W_m=0$ 更强。
而\textbf{人工绝缘球域并非无限导电海水}。
改变 $R$ 还改变全域包络；
外球对照不能在没有固定试探空间的情况下
认作单纯截断误差。
可通过完整内外区非局部边界算子或边界积分来
模拟无限外域；相应理论与 FEM--BEM 对称耦合见
Hiptmair~\cite{hiptmair2001coupling}。
局部球谐外推已实现，但不能充当双盒内外
连续场与边界迹双向匹配的证据。

\subsection{Coulomb 磁能的核心问题是奇异积分的一致性}
目标仍是\emph{原核}：
\begin{equation}
L_{mn}=\frac{\mu_0}{4\pi}
\int_\Omega\int_\Omega
\frac{\bm W_m(\bm x)\cdot\bm W_n(\bm y)}
{|\bm x-\bm y|}\,\mathrm d\bm y\,\mathrm d\bm x.
\label{eq:g1-review-original-coulomb}
\end{equation}
任意求积权重所造的等体积球形电流云，
即使获得半正定的磁能 Gram 矩阵，
仍不是式 \eqref{eq:g1-review-original-coulomb}
的一致性证明。

针对涡流体积分方法的直接先行文献
Passarotto--Pitassi--Specogna~\cite{passarotto2022}
对电感双积分的自项、相邻项及电流空间构造
开展专门分析，并指出简单采用两组内外数值
求积可使磁矩阵对角积分出现严重错误。
这与本项目 $x/y$ 分轴加密大幅非加性相呼应；
不能以双轴变化较小来认定两种错误相消后的答案正确。
奇异抽取与 Duffy 型变量替换参见
\cite{jarvenpaa2003,duffy1982}。

\paragraph{优先实验的等价连续奇异减除公式}
对固定点 $\bm x$、任何满足所需局部正则性的
模式 $\bm W_n$ 和局部径向窗 $\chi_a$，
可写严格恒等式
\begin{align}
\int_\Omega\frac{\bm W_n(\bm y)}{|\bm x-\bm y|}
\,\mathrm d\bm y
={}&\int_\Omega
\frac{\bm W_n(\bm y)-\chi_a(|\bm x-\bm y|)\bm W_n(\bm x)}
{|\bm x-\bm y|}\,\mathrm d\bm y\nonumber\\
&+\bm W_n(\bm x)
\int_{\Omega\cap B_a(\bm x)}
\frac{\chi_a(|\bm x-\bm y|)}{|\bm x-\bm y|}
\,\mathrm d\bm y.
\label{eq:g1-review-subtraction}
\end{align}
第一项在近对角的主导常值奇异性被抵消，
第二项是仅与真实边界几何有关的标量积分；
可利用物理球坐标 $\rho^2\,\mathrm d\rho/
\rho=\rho\,\mathrm d\rho$ 和旋转盒精确射线裁剪。
公式本身不引入核软化，但仍要求
对盒角、边缘、切触方向和场低正则性进行
\emph{独立的局部积分误差控制}。
优先以解析可核的固定小模态测试该式，
再研究原核 $\bm L$ 的对角/相邻元素，
不立即扩张至 229 模态。

\subsection{端口两种等价弱式是积分误差探针}
对 $\bm W_m=\Curl\bm F_m$，
由积分恒等式
\[
\int_\Omega(\Curl\bm F_m)\cdot\bm C_p
=\int_\Omega\bm F_m\cdot(\Curl\bm C_p)
+\int_{\partial\Omega}(\bm n\times\bm F_m)
\cdot\bm C_p\,\mathrm dS.
\]
若整个边界 $\bm n\times\bm F_m=0$ 且
$\bm B_{\rm src,p}=\mu_0\Curl\bm C_p$，则
\begin{equation}
\boxed{K_{mp}
=\mu_0\int_\Omega \bm W_m\cdot\bm C_p\,\mathrm dV
=\int_\Omega\bm F_m\cdot\bm B_{\rm src,p}\,\mathrm dV.}
\label{eq:g1-review-kdual}
\end{equation}
现有离散线源的 $\bm B$ 正是其节点势核 $\bm C$
的解析旋度；故离散的两种 $K$ 公式相差
百分之几至几十个百分点，
主要反映体积分欠分辨（或基/源码调用不一致），
不是两个不同的物理理论。
将 sphere 与 shell 分开算时，
二者界面可能产生额外边界项，
不能要求每个子分区各自满足 \eqref{eq:g1-review-kdual}。

\subsection{正权、QMC 与离散能量恒等式的正确用途}
正权求积能保正的电阻 Gram 矩阵；
利用球形电流云或别的光滑核也可能保正的
\emph{替代}磁能。但正定、互易与被动性
不能取代真正的原核积分误差分析。
扰乱 Sobol QMC 可以对 $\bm R,\bm K$
提供独立构造的交叉检查，
但靠单次随机扰乱差异不能得到严格
连续积分的高精度证书。
非光滑/不连续被积函数的随机 QMC 收敛
另有条件与速率分析~\cite{hewang2015}；
复杂边界层不能直接套用光滑被积函数的最优阶。

\subsection{严格分关口的参考验证合同}
\begin{enumerate}
\item \textbf{源和物理域：} 锁定 BFZI 连续线圈、
闭合回流、源节点阶次、盒壁和材料参数，
确认电流源连续性、C/B 同源旋度及 MQS 适用尺度。
\item \textbf{原核积分：}
在解析可解几何和固定电流场中，
分别控制原 Coulomb $\bm L$ 的对角、
相邻、远区块；与独立奇异抽取和另一构造比较。
不能用 uniform-ball 的半正定代替一致性。
\item \textbf{R/K：}
同一连续模式下独立计算每个矩阵元素，
在盒面/边/角与球远区分开定量估计；
完整域上式 \eqref{eq:g1-review-kdual}
须随加密收敛到同一极限。
\item \textbf{空间与外域：}
在积分可靠之后检验真正嵌套的函数空间和
固定边界算子；比较合规无限域续接或
独立外部物理参考，区分外球误差与模式误差。
\item \textbf{场和代理：}
使用固定独立参考检验 $\bm Z$、
海水 $\bm J$、盒壁/角部焦耳热及总损耗；
再建立热场和跨几何神经教师，必要时允许
\emph{离线} FEM/BEM 作裁判，不迫使在线代理使用固定网格。
\end{enumerate}

TEAM 7 的跨方法与实测基准说明需要
独立外部裁判而非单实现自洽性~\cite{fujiwara1990team7}，
但 TEAM 7 的有孔导电体不是当前双实心绝缘盒
外部导电海水问题，不能替代 BFZI 自身基准。
上述关口完成以前，
\textbf{G1 仍不通过；本节不批准新的正式神经教师。}

% BFZI P0, 2026-11-04. Source-only derivation and validation contract.
\section{P0：锁定磁准静态物理域、连续闭合源项与有限直导线参考}
\label{sec:p0-physical-source}

\subsection{单一物理场景与相量符号}
在本节首先只讨论静止、线性、$\mu=\mu_0$、
海水标量电导率 $\sigma_s>0$、
两只互不相交的理想绝缘实心封装及
BFZI 圆形--圆角方形\textbf{单根连续螺旋}。
每个端口的电流沿螺旋外端到内端，
再沿显式规定的封装内回流折线回到外端，
形成相同方向约定下的闭合曲线 $\Gamma_p$。
正频相量为 $e^{+\jj\omega t}$，端口电流使用峰值。
圆形情形是角部连续绕行与直边长度为零的
合法极限，不重新引入旧超椭圆族。

海水外部真实区域与人工有限计算区域不同：
\[
\Omega_\infty=\RR^3\setminus(\overline B_1\cup\overline B_2),
\qquad
\Omega_R=B_R\setminus(\overline B_1\cup\overline B_2).
\]
封装界面满足 $\bm n\cdot\bm J_s=0$，
而计算球的相同零法向条件只是
\emph{人为的有限域电流条件}。
在未对人工球和无限海水外区实现非局部相容耦合前，
不得称 $\Omega_R$ 的离散解为连续无限域真值。

\subsection{涡流近似的必要尺度核查}
海水中忽略位移电流并不等于自动满足
Maxwell--MQS 的所有适用条件：
至少需检验
\[
\eta_{\mathrm D}
=\frac{\omega\epsilon_s}{\sigma_s}\ll1,
\qquad
\eta_{\mathrm wave}
=k_0D=\frac{\omega D}{c_0}\ll1,
\qquad
\delta_s=\sqrt{\frac{2}{\omega\mu_0\sigma_s}} .
\]
还要检查真实介质波数与装置尺度、
源电流的连续性、界面极化/介电损耗是否可忽略。
Ammari--Buffa--N\'ed\'elec~\cite{ammaribuffanedelec2000}
强调涡流近似的高阶准确性依赖
外加源所满足的额外兼容条件。
例如 $f=200\,\mathrm{kHz}$、
$\sigma_s=5\,\mathrm{S/m}$、
$D=0.52\,\mathrm m$ 时
$\delta_s\approx0.503\,\mathrm m$，
$k_0D\approx2.18\times10^{-3}$；
\emph{仅作例示}若设海水相对介电常数为 $80$，
则 $\eta_{\mathrm D}\approx1.78\times10^{-4}$。
当前海水模型尚无实际的色散介电输入，
这些示例数字\textbf{不提供任意海水工况的电磁误差界}。

\subsection{同一闭合路径的理想线源核}
将外部指定单位电流源记作
\[
\bm C_p(\bm x)=\frac1{4\pi}
\oint_{\Gamma_p}\frac{\mathrm d\bm\ell'}{|\bm x-\bm r'|},
\qquad
\bm A_{\mathrm{src},p}=\mu_0\bm C_p,\qquad
\bm B_{\mathrm{src},p}=\mu_0\Curl\bm C_p .
\]
它只在实际海水域中评价，避免零半径导线上的
奇异核。闭合曲线积分与
\emph{有限节点点线元加权和}不是同一个分布：
\[
\bm C_{p,h}(\bm x)=\frac1{4\pi}
\sum_k\frac{\bm h_k}{|\bm x-\bm r_k|}.
\]
两者在光滑、离线距离有正下界的查询区域可通过
积分阶次收敛相互逼近。
虽然 $\bm B_{p,h}=\mu_0\Curl\bm C_{p,h}$
对上述同一核严格成立，
且 $\sum_k\bm h_k\approx0$，
仍不自动使有限点线元分布成为
严格无散度的物理闭合回路电流。
必须独立检查盒壁附近的
$\bm C_{p,h}$、$\bm B_{p,h}$ 求积误差。

\subsection{不依赖源 Gauss 节点的精确有限直导线核}
对起止点 $\bm a\to\bm b$，
令 $L=|\bm b-\bm a|$、
$\bm t=(\bm b-\bm a)/L$、
$\bm R_1=\bm x-\bm a$、
$\bm R_2=\bm x-\bm b$、
$r_1=|\bm R_1|$、$r_2=|\bm R_2|$。
从线积分直接计算得到
\[
\boxed{
\bm C_{\mathrm{seg}}(\bm x)
=\frac{\bm t}{4\pi}
\log\frac{r_1+r_2+L}{r_1+r_2-L}
}
\]
和
\[
\boxed{
\bm B_{\mathrm{seg}}(\bm x)
=\frac{\mu_0}{4\pi}
\frac{(\bm R_1\times\bm R_2)(r_1+r_2)}
{r_1r_2(r_1r_2+\bm R_1\cdot\bm R_2)} .
}
\]
仅在不位于理想导线的场点使用。
这两个公式无近奇异 Gauss 点问题，
可以作为 BFZI 每根\textbf{真实直边}
及每段明确回流线的独立源场参考。
BFZI 角部仍然是连续曲线：
固定角部几何和真实有向线元，
仅对角部使用高阶 Gauss 收敛，不得
用圆弧中心拟合、独立闭合匝或假定匝同心
来替代实际 BFZI 公式。

代码 \texttt{sdfmpneo\_bfzi/source\_reference.py}
提供 \texttt{straight\_segment\_field} 与
\texttt{continuous\_port\_field}。
首轮 P0 本机 \texttt{a50ed62} 的
75 项测试通过；独立角段 24/32 阶积分一致至机器精度，
但旧生产源的 4 阶整路径 Gauss
$\bm C/\bm B$ 探针误差分别约
$0.253\%$ 和 $1.174\%$，
12 阶时分别约 $0.000118\%$ 和 $0.00380\%$。
因此旧 4 阶源场\textbf{未达到}预先规定的
$10^{-3}$ 门槛。

随后已经据此修正\texttt{LineSource}：
实际端口源场 $\bm C$ 与 $\bm B$ 的直边
和全部回流段使用上述解析有限直导线核，
曲角才采用该阶次的 Gauss 节点。
完整的旧线元节点仍独立保留用于导线长度、
铜线热源和空气 Neumann 互感近似，
避免错误地把上述源场修正称为已修复铜线
邻近效应或完整空气自感。
源场两种读出继续使用同一
$\bm C$ 的解析旋度，即
$\bm B=\mu_0\Curl\bm C$；
本机在代码 \texttt{ce086b1} 上最新
77 项测试通过，并在两端口各 28 个海水探针上测得：
新 4 阶 $\bm C/\bm B$ 的相对误差分别为
$2.24\times10^{-5}$ 与 $1.95\times10^{-4}$，
原全 Gauss 4 阶则为
$2.53\times10^{-3}$ 与 $1.17\times10^{-2}$。
故\textbf{固定示例上的源场已通过预定 $10^{-3}$ 门槛}，
不能将该点集结果外推全部几何参数域或
真实有限截面导线近邻电磁分布。

\paragraph{闭合回路的 Coulomb 规范连续性检查}
完整的光滑闭合源曲线在导线外满足
\[
\Div\bm C_p(\bm x)
=-\frac1{4\pi}\oint_{\Gamma_p}
\frac{\mathrm d\bm\ell'\cdot(\bm x-\bm r')}
{|\bm x-\bm r'|^3}
=-\frac1{4\pi}\oint_{\Gamma_p}
\mathrm d_{\ell'}\frac1{|\bm x-\bm r'|}
=0.
\]
对任何直线段 $\bm a\to\bm b$，
这项的精确贡献为
\[
\Div\bm C_{\rm seg}(\bm x)=
\frac1{4\pi}\left(
\frac1{|\bm x-\bm a|}
-\frac1{|\bm x-\bm b|}\right).
\]
沿物理闭合路径各段端点项望远镜式抵消；
仅曲角用有限 Gauss 时残留可测数值偏差。
现在实现的
\texttt{LineSource.divergence\_kernel}
直接积分这些解析端点项和曲角点线元的导数核，
而不是通过 $B=\mu_0\Curl C$ 自我验证。
同时新增盒面、棱、角外点的局部源场误差。
本项扩展诊断尚待本机运行，不作为新的
全 Maxwell 认证或非磁性有限截面适用域证明。


\subsection{P0 连续 BFZI 几何的角部正向性约束}
支撑函数角部的有向切线在两种绕向下都满足
\[
\frac{\mathrm d\bm r_{m,k}}{\mathrm du}
=\sigma_w\!\left[\varrho(\eta)-
2\sigma_w(1-\alpha)\beta
(\bm e_m\cdot\bm t(\psi_m))\right]\bm t(\psi_m),
\qquad
\sigma_w(\bm e_m\cdot\bm t)=\cos u-\sin u\in[-1,1].
\]
所以，\emph{只检查 $\varrho(\eta)>0$ 不足以防止角部
切线消失或反向}。
目前的 BFZI 构造器加上一个直接有理论依据的
保守充分条件（并非必要条件）：
\[
\boxed{\;
\varrho_{\min}=\varrho_0-\alpha sN
>2(1-\alpha)\beta,\qquad\beta=s/(2\pi).\;}
\]
若违反该条件，
不得把仍为正数的角半径误当成一根正向无折返的
圆角方形连续导线。这避免错误源几何进入求积和
后续海水教师。有限半径导线的自接触/弯曲半径
尚需更强的几何条件，不能从角切线正向性单独推出。

\subsection{严格闭合折线参考与连续曲角近似误差}
本机 \texttt{325d47b} 的 P0 日志（79 项测试通过）
将双端口的查询从 28 点扩至各 84 点，
增加真实盒面/棱/角附近的海水点：
4 阶源的 $\bm C,\bm B$ 总体误差分别
$7.28\times10^{-5},4.27\times10^{-4}$，
边角 $\bm B$ 误差 $7.48\times10^{-4}$；
而该源的归一化 $\Div\bm C$ 仍为
$1.06\times10^{-4}$，从 6、8、12 阶
分别下降至
$3.31\times10^{-6},1.02\times10^{-7},8.86\times10^{-11}$。
这个结果支持曲角 Gauss 残差逐级收敛，
\emph{但它不是严格满足无散度源条件的离散源}。

因此另行构造真正闭合的折线
$\bm v_0,\ldots,\bm v_M=\bm v_0$，
每段为从解析 BFZI 曲角上选取的相邻真实位置
形成的弦、完整解析直边或原回流折线。
所有段仍使用未经软化的有限直线精确核；
任意离散阶数、任意导线外查询位置，
其散度严格望远镜式抵消：
\[
\Div\bm C_{\rm poly}(\bm x)
=\frac1{4\pi}\sum_{j=0}^{M-1}
\left(\frac1{|\bm x-\bm v_j|}
-\frac1{|\bm x-\bm v_{j+1}|}\right)
=0.
\]
该性质\textbf{完全独立于曲角 Gauss 求积}；
且曲角只在正切线的原 BFZI 路径上细分，
所以不引入任何历史超椭圆或独立闭合同心圆。
然而折线与实际连续曲线并不相同，
要使 $\bm C_{\rm poly},\bm B_{\rm poly}$
逼近连续源仍必须做角部弦段
16/32/64 个面板的\emph{几何收敛}比较。
闭合性是物理合法性，场精度是另一个量。
代码接口为
\texttt{polygonal\_port\_vertices} 和
\texttt{polygonal\_port\_field}。
其新增实验与测试尚待本机验证；
暂不切换默认生产源算法。

\subsection{P0 验收和与 P1 的严格分界}
P0 必须在固定两端口位姿、
真实连续中心线和完整回流下执行：
\begin{enumerate}
\item 直边、角部及相邻匝真实首尾连接；
线圈外端到内端再沿回流回到外端；
旋转变换和源向量方向符合右手坐标。
\item 在与真实直线源远离的场点，解析
$\bm C_{\mathrm{seg}},\bm B_{\mathrm{seg}}$
与足够高阶的独立 Gauss 源积分相等；
检查刚体协变和
$\bm B_{\mathrm{src}}=\mu_0\Curl\bm C_{\mathrm{src}}$。
\item 圆形与圆角方形均以独立参考阶次
\texttt{corner\_order=24/32} 比较场；
固定盒壁、边缘、角点和远场查询点，
报告生产 \texttt{order=3/4/6/8/12}
的源场相对误差，不以端口能量闭合替代。
\item 对每个目标介质/频率重新核查
$\omega\epsilon_s/\sigma_s$、$k_0D$、
皮肤深度和内部导体材料近似。
本节不认证超出线性均匀海水的情形。
\item P0 的正确闭合源只是
下一阶段原 Coulomb 磁电感积分的
\emph{已冻结外部激励}；在磁能
$L_{mn}$ 尚未独立收敛前，
不能宣称海水 $J$、$q$、阻抗或神经教师达到真值。
\end{enumerate}
本机入口是
\texttt{python run.py --mode bfzi-p0-source --with-tests}，
将只输出源项与尺度的
\texttt{p0\_physics/p0\_reference/p0\_source}，
不启用 GitHub Actions。
原解析源修正已在本机 77 项测试中通过，新增闭合\texttt{div C}
与盒棱/角源场扩展诊断尚待本机复测。
当前\textbf{P0 固定示例源精度通过，但参数域尚未全部验收，G1=FAIL}。

% 2026-11-04. P0 limited source acceptance and exact P1 Coulomb test.
\section{P0 限定源项冻结与 P1 原 Coulomb 磁能的解析关口}
\label{sec:p1-exact-coulomb-gate}

\subsection{P0 的可验证结论及适用域}
本机在 \texttt{21daadf} 上
83 项单元测试通过，P0 两端口各 84 个海水
面/棱/角/远场采样点上，
生产源四阶解析直线与曲角 Gauss 方法
的 $\bm C$、$\bm B$ 相对误差为
$7.28\times10^{-5}$、$4.27\times10^{-4}$，
曲角 24--32 阶高阶独立参考近机器一致。
另一种基于原 BFZI 曲角的物理闭合折线
对 16、32、64 弦段的相对磁场误差
分别约 $3.90\times10^{-4}$、
$9.76\times10^{-5}$、$2.44\times10^{-5}$，
呈近似二阶衰减，且每一折线
$\Div\bm C_{\rm poly}$ 在导线外均为机器零。

因此可以将\emph{给定物理参数和理想闭合单位线源}
的 P0 源积分定义与对照视为限定通过，
不能将之拓展为全参数域、有限线径的自碰撞
或真实铜线近邻电流分布的认证。
固定相量方向、回流路径和材料条件，
并明确人工有限绝缘外球仍非无限导电海水。

\subsection{严格符合电流约束的解析球形原核基准}
在 $\Omega_R=\{\bm x\in\RR^3:|\bm x|<R\}$，
选取三个真实电流向量基：
\[
\bm W_i(\bm x)=\bm e_i\times\bm x,\qquad i=1,2,3.
\]
各模式满足
\[
\Div\bm W_i=0,\qquad
\bm n\cdot\bm W_i|_{\partial\Omega_R}=0,
\qquad
\bm W_i=\Curl\bm F_i,\qquad
\bm F_i=\frac{R^2-|\bm x|^2}{2}\bm e_i,
\quad
\bm F_i|_{\partial\Omega_R}=0 .
\]
这是具有连续物理零法向边界条件和零迹旋度势的
严格电流模式，不是随意选择的常向量试验。

对任一笛卡尔分量为 $\bm W_i$ 的线性函数，
定义原 Coulomb Newton 势
\[
\bm U_i(\bm x)
=\int_{\Omega_R}
\frac{\bm W_i(\bm y)}{|\bm x-\bm y|}\,
\mathrm d\bm y .
\]
利用 $\Delta_{\bm x}\bm U_i=-4\pi\bm W_i$，
球内 $l=1$ 正则径向解
$\bm U_i=c(r)\bm W_i$，
以及球外衰减偶极项的连续匹配，可得
\begin{equation}
\boxed{
\bm U_i(\bm x)
=\left(\frac{2\pi R^2}{3}
-\frac{2\pi|\bm x|^2}{5}\right)
\bm W_i(\bm x),\quad |\bm x|\le R .
}
\label{eq:p1-ball-exact-potential}
\end{equation}
这里不含任何软化、替代球云电流模型或数值
对角自项定义，且边界与标准 $1/r$ 核均为真实积分。

球的各向同性使不同坐标轴模式正交。
由于
\[
\int_{\Omega_R}|\bm W_i|^2\,\mathrm d\bm x
=\frac{8\pi R^5}{15},\qquad
\int_{\Omega_R}|\bm x|^2|\bm W_i|^2
\,\mathrm d\bm x=\frac{8\pi R^7}{21},
\]
代入原核双重积分直接得到
\begin{align}
L_{ij}
&=\frac{\mu_0}{4\pi}
\iint_{\Omega_R\times\Omega_R}
\frac{\bm W_i(\bm x)\cdot\bm W_j(\bm y)}
{|\bm x-\bm y|}
\,\mathrm d\bm y\,\mathrm d\bm x\\
&=\boxed{\frac{16\pi\mu_0 R^7}{315}\,\delta_{ij}}.
\label{eq:p1-ball-analytic-L}
\end{align}
这些模式\emph{没有按安培载流量归一化}，
故其算子双线性型带模式尺寸因子，不得
把该数值误标为真实两端口电感的亨利数值。
本关以完全已知的对角/非对角三模态磁能矩阵
提供\textbf{可证伪的原核数值参考}。

\subsection{相同电流、不同积分轴与局部核分区}
\texttt{examples/bfzi\_p1\_coulomb.py}
使用独立球形 Gauss \texttt{x/y} 积分规则，
直接调用当前未软化原核
\texttt{assemble\_coulomb\_split}：
\begin{enumerate}
\item 固定上述 $\Omega_R,\bm W_i$，
将外积分轴 $x$、内积分轴 $y$ 以
\texttt{base/x\_only/y\_only/both} 四组比较
\emph{同一个解析真值}；
\item 固定节点和电流，改变局部
分区球半径 $a/R=0.25,0.55$，
检查结果对\emph{数学上等价的原核分区}
的敏感度，不改变核 $1/r$；
\item \textbf{在本机首次运行前冻结球基准专用验收阈值：}
用最终 \texttt{both} 规则与式
\eqref{eq:p1-ball-analytic-L} 比较，
矩阵 Frobenius 相对误差、
最大对角元相对误差、非对角相对误差
均须小于 $10^{-3}$；
以解析 $L_{11}$ 归一化的
局部分区 $a/R=0.25,0.55$
双规则差也须小于 $10^{-3}$。
\texttt{base/x\_only/y\_only/both} 的序列仍须人工
分析是否稳定逼近同一真值，
不允许只因最后一个样本偶然落在阈值内就认证。
\item 记录原双重积分三模态矩阵与式
\eqref{eq:p1-ball-analytic-L} 的整体
Frobenius 相对误差、最大对角相对误差、
非对角误差及对称性偏差；
\item 若球基准未达到有证据的稳定性，
先查核奇异区、局部求积与
内外积分轴，而不继续扩大 229 模态字典，
不更换成 Plummer 或实体球云近似。
\end{enumerate}
本机独立入口为
\texttt{python run.py --mode bfzi-p1-coulomb --with-tests}，
独立版本化日志在
\texttt{docs/research/bfzi\_uwpt\_electrothermal/runs/p1/}。
当前\textbf{该 P1 新算例尚未取得本机数值运行证据}，
不能宣布原核求积成功。


\subsection{P1 原核球基准实测失败及严格误差归因}
用户本机 \texttt{ba16d6a} 的第一个 P1 运行
85 项程序单测通过，但\emph{原核解析准确度失败}。
在 $R=1\,{\rm m}$ 的三解析无散度模式上，
当前分区算子 \texttt{assemble\_coulomb\_split}
分别以基础、只加密 $x$、只加密 $y$、
双轴加密的体积规则得到 Frobenius 相对误差
$2.120667\%$、$1.556394\%$、
$1.146122\%$、$0.614486\%$，
最后一项仍高于预定 $0.1\%$ 目标。
固定 $x/y$ 规则，局部分区半径 $a/R$
改为 $0.25,0.55$ 后，
$L_{11}$ 的差值相对解析真值为 $0.4230034\%$。
对称性、非对角元误差均接近机器零，
没有提供数值正确性的充分证据。

目前的 C$^2$ 局部径向窗函数满足
$\chi(s)=1$（$s\le1/2$），
$\chi(s)=0$（$s\ge1$），中间为
\[
\chi(s)=1-10(2s-1)^3+15(2s-1)^4-6(2s-1)^5,
\qquad \frac12<s<1,
\]
故其局部矩可\emph{完全解析}：
\[
\int_0^1 s\chi(s)\,\mathrm ds=\frac27 .
\]
令 $a(\bm x)=
\min(a_{\max},0.6\,\mathrm{dist}(\bm x,\partial\Omega_R))$，
局部球包含于导电球内部，
且 $\bm W_i(\bm y)=
\bm W_i(\bm x)+\bm e_i\times(\bm y-\bm x)$；
后项在对称球上积分为零，于是
\begin{equation}
\boxed{
\int_{|\bm y-\bm x|<a(\bm x)}
\frac{\chi(|\bm y-\bm x|/a)}
{|\bm x-\bm y|}\bm W_i(\bm y)\,\mathrm d\bm y
=\frac{8\pi}{7}\,a(\bm x)^2\,\bm W_i(\bm x).
}
\label{eq:p1-analytic-local-window}
\end{equation}
这一式直接给出\emph{原积分分区}的确切局部值，
无需再加密局部球的径向、极向或方位求积。
已有球内势的解析解
\eqref{eq:p1-ball-exact-potential} 则可直接在
$x$ 规则上计算原算子而不经过任何 $y$ 求积。
两者相减就能将误差定位为：
\begin{enumerate}
\item 解析内势经 $x$ 求积的外积分误差；
\item 局部球采用式
\eqref{eq:p1-analytic-local-window} 精确值后，
$y$ 远区规则相对解析内势的误差；
\item 原程序给出的局部球数值值与上述精确值的差。
\end{enumerate}
其中远区核
$(1-\chi(|\bm x-\bm y|/a(\bm x)))/|\bm x-\bm y|$
虽然在对角线取零，但靠近球壁时
$a(\bm x)\to0$，转折层收缩。
固定全局球 Gauss 求积容易漏掉这类细小空间尺度。
新 \texttt{p1\_decomp} 旨在直接证实或反驳
该误差来源，不能凭直觉认定它已经通过。

\subsection{保留物理原核的全域奇异常量减除}
对任何具有足够正则性的测试场，恒等式
\begin{equation}
\boxed{
\int_\Omega\frac{\bm W(\bm y)}{|\bm x-\bm y|}
\,\mathrm d\bm y
=\bm W(\bm x)\,S_\Omega(\bm x)
+\int_\Omega
\frac{\bm W(\bm y)-\bm W(\bm x)}{|\bm x-\bm y|}
\,\mathrm d\bm y,
\quad
S_\Omega(\bm x)=\int_\Omega
\frac{\mathrm d\bm y}{|\bm x-\bm y|}.
}
\label{eq:p1-global-constant-subtraction}
\end{equation}
当 $\bm W$ 局部 Lipschitz 时，
最后一项在 $y\to x$ 的函数值有界，
但一般仍有依赖方向的极限，不得在
$x/y$ 节点重合处任意赋值。
这是\textbf{未软化的原始 Coulomb 核}，
且并非均匀小球电流云替代模型。

球形几何的标量势已知：
\[
S_{B_R}(\bm x)=2\pi
\left(R^2-\frac{|\bm x|^2}{3}\right).
\]
据此新增
\texttt{green\_coulomb.assemble\_coulomb\_subtracted}，
直接用给定真实几何的独立标量势
和独立 $x/y$ 体积节点，对式
\eqref{eq:p1-global-constant-subtraction}
构成原核磁能矩阵，作为第二种 P1 算法候选。
当前仅能在\emph{球上}把解析
$S_{B_R}$ 当成数值真值。

对双盒排除的海水域 \emph{不可以直接使用}
球形 $S_{B_R}$。
若要使该算法应用于真实
$\Omega=B_R\setminus(\overline B_1\cup\overline B_2)$，
仍需独立构造并验证
$S_\Omega(\bm x)$。
用散度恒等式
$\Div_{\bm y}[(\bm y-\bm x)/|\bm y-\bm x|]
=2/|\bm y-\bm x|$，
可将其形式上写成物理边界的曲面积分：
\[
\boxed{
S_\Omega(\bm x)=\frac12
\int_{\partial\Omega}
\frac{\bm n_{\bm y}\cdot(\bm y-\bm x)}
{|\bm y-\bm x|}\,\mathrm dS_{\bm y}.
}
\]
外球使用外法向，而内盒壁使用\emph{海水域外法向}
（即朝向绝缘盒内部）。
该表面表示本身需要近壁误差控制，
\textbf{在未经证明或核对前不能称双盒算法已通过}。
已新增诊断和原核减除实现，
但这些变更尚未取得用户本机 P1 新实测。


\subsection{原核进一步减除一阶几何矩：非线性解析涡流验算}
仅使用式 \eqref{eq:p1-global-constant-subtraction}
减去常量时，剩余被积核虽然有界，
但对角附近仍具有角向极限。
对有合适 $C^2$ 正则性的电流，
可以进一步将 \emph{完全相同的物理原核}
分解为
\begin{align}
\bm U_W(\bm x)
={}&\bm W(\bm x)\,S_\Omega(\bm x)
+\nabla\bm W(\bm x)\,\bm M_\Omega(\bm x)
\notag\\
&+
\int_\Omega
\frac{
 \bm W(\bm y)-\bm W(\bm x)
 -\nabla\bm W(\bm x)(\bm y-\bm x)}
 {|\bm x-\bm y|}\,\mathrm d\bm y ,
\label{eq:p1-affine-taylor-subtraction}\\
\bm M_\Omega(\bm x)
={}&\int_\Omega\frac{\bm y-\bm x}{|\bm x-\bm y|}
\,\mathrm d\bm y
=\int_{\partial\Omega}
 |\bm x-\bm y|\bm n_{\bm y}\,\mathrm dS_{\bm y}.
\label{eq:p1-geometry-first-moment}
\end{align}
最后一式由
$\nabla_{\bm y}|\bm x-\bm y|=
(\bm y-\bm x)/|\bm y-\bm x|$
及散度定理给出。
余项因一阶 Taylor 减除在对角仅为
$O(|\bm x-\bm y|)$，
从数学上比原始 $1/r$ 核和零阶减除更规则。
\textbf{式 \eqref{eq:p1-affine-taylor-subtraction}
是严格恒等式，不是软化或新的物理模型}。

球内的两个几何矩具有独立的解析值：
\[
S_{B_R}(\bm x)=
2\pi\left(R^2-\frac{|\bm x|^2}{3}\right),\qquad
\boxed{\bm M_{B_R}(\bm x)=
\frac{4\pi}{15}(|\bm x|^2-5R^2)\bm x.}
\]
若仅用一阶电流 $\bm W_i=\bm e_i\times\bm x$，
Taylor 余项\emph{严格为零}，因而该例只能
验证 $\bm S,\bm M$ 和球外积分，
\textbf{不能}据此认证余项的实际求积。

为防止上述制造解过于平凡，
新增一个\emph{非线性、严格无散度、
球壁零法向} 的第二级真值模式：
\[
\bm W_i^{(3)}(\bm x)
=|\bm x|^2(\bm e_i\times\bm x),\qquad
\bm F_i^{(3)}(\bm x)
=\frac{R^4-|\bm x|^4}{4}\bm e_i ,
\quad \Curl\bm F_i^{(3)}=\bm W_i^{(3)} .
\]
从 $\Delta\bm U_i^{(3)}=-4\pi\bm W_i^{(3)}$
及球外 $l=1$ 匹配得到
\begin{equation}
\boxed{\bm U_i^{(3)}(\bm x)
=\left(\frac{\pi R^4}{3}
-\frac{\pi|\bm x|^4}{7}\right)
(\bm e_i\times\bm x),\quad
L_{ij}^{(3)}=\frac{16\pi\mu_0 R^{11}}{693}\delta_{ij}.}
\label{eq:p1-cubic-vortex-exact}
\end{equation}
这个模式的 Taylor 余项一般\emph{不为零}，
可以真正检验式
\eqref{eq:p1-affine-taylor-subtraction}
的残余 $y$ 积分。
独立短程球体积数值计算在固定 480 个 $x$、
756 个 $y$ 节点下发现：
仅减常量的非线性电流磁矩阵相对误差约
$2.453\times10^{-3}$，
进一步减除一阶项后约
$1.596\times10^{-4}$，
而解析内势仅在 $x$ 节点积分类似误差约
$5.91\times10^{-6}$。
这些是独立短程校验数值，
\textbf{新增 GitHub P1 自动日志与单测尚待本机复跑}，
不能将其等同于完整 BFZI 教师验收。

生产代码的\texttt{assemble\_coulomb\_subtracted}
现在支持配对参数
\texttt{gradient}、\texttt{first\_moment}，
并坚持输入 \emph{已独立可信的}
实际几何标量势和一阶矩。
对有两个旋转绝缘盒孔洞的海水域，
不能直接代入球公式：
\textbf{真实双盒 $S_\Omega,\bm M_\Omega$
和临近盒壁的相应积分必须分别认证}。
这一步没有改动原 BFZI 端口源、
$R/K$ 或神经网络教师训练。

\subsection{与真实 BFZI 双盒积分之间的限制}
即使式 \eqref{eq:p1-ball-analytic-L} 在球上成功，
它只证明在此解析无散度电流空间和球形边界
下的原 Coulomb 核积分经过了一个独立基准。
BFZI 的真实海水域是
两只独立 SE(3) 旋转绝缘盒的外部，
存在非光滑棱角和多体近邻相互作用，
另需对\emph{同一固定物理电流场} 的对角块、
靠近盒面的局部积分与远区子块单独认证。
涡流体积积分中的奇异电感自项与
分离内外积分轴的困难见
Passarotto--Pitassi--Specogna~\cite{passarotto2022}；
式 \eqref{eq:g1-review-subtraction} 的解析
奇异减除仍是后续优先途径，
不可以用球形算例代替所有边界角度验证。


% Frozen physical direct problem, 2026-11-06; no solver accuracy certified.
\section{锁定 BFZI 物理直接法：Maxwell--MQS、独立参考和验收}
\label{sec:bfzi-physical-direct-locked}

\paragraph{锁定决定。}
此节以及
\texttt{docs/research/bfzi\_uwpt\_electrothermal/PHYSICAL\_DIRECT\_METHOD\_LOCKED.md}
和
\texttt{PHYSICAL\_DIRECT\_METHOD\_EVIDENCE\_AND\_GATES.md}
是物理直接法 V1.0 的优先约束。历史 P1-C/G1
真实海水五模态原核积分仍未收敛，保持 FAIL。
不要把对称、正权积分的代数性质或两种相近的角向规则
当作物理连续解精度的证明。正式神经教师保持暂停。

\subsection{物理任务与文献适用性}
用户提供的 2026-11-05 水下 WPT 综述严格区分
A 类代理、B 类前向物理求解、C 类在线参数辨识。
当前先完成 B 类。
Shafiei 等~\cite{shafiei2025uwpt}
在共轴圆形平面线圈与分层导电介质的
特定 Helmholtz 模型下计算海水附加电阻和互阻；
Wang 等~\cite{wang2025marinewpt}
将复自感和复互感与场路联立；
Fei 等~\cite{fei2025seawater}
通过提取模式和拟合评估互阻及相位影响。
这些结果不是任意 SE(3) 双盒和连续圆角方形
BFZI 绕线的通用闭式真值。
成熟三维独立参考采用合规的 $A$--$\phi$ 涡流法，
必要时配合 FEM--BEM 开域耦合
~\cite{hiptmair2001coupling,hiptmaiostrowski2005}；
仅在材料和边界假设匹配时考虑 BEM
~\cite{tsuboitanaka1990}。

\subsection{固定真实场景和源}
两独立、互不穿透的非导电、非磁性实体绝缘盒
$B_1,B_2$ 允许各自变尺寸及刚体 SE(3) 位姿；
每盒包含圆形或圆角方形\emph{单根连续多匝 BFZI}
线圈及明确的完整闭合馈回流，不是超椭圆。
海水基线：
\[
 \Omega_\infty=\RR^3\setminus
 (\overline{B_1}\cup\overline{B_2}),
 \qquad \sigma_s>0,\quad\mu=\mu_0.
\]
两个端口峰值相量 $\bm I=(I_1,I_2)^T$
取 $e^{+\jj\omega t}$。
源须满足分布连续性；有限直径铜的自阻抗
不能由零半径闭合线源的发散自积分代替。
磁准静态近似必须先检查频率、尺度、
电导率和位移电流以及源兼容性
~\cite{ammaribuffanedelec2000}。

\subsection{锁定 Maxwell--MQS 与 $A$--$\phi$}
\begin{align}
 \Curl\bm H&=\bm J_{\rm src}+\bm J_e,&
 \Curl\bm E&=-\jj\omega\bm B,\\
 \Div\bm B&=0,&
 \bm B&=\mu\bm H,&
 \bm J_e&=\sigma\bm E.
\end{align}
海水域 $\Div\bm J_e=0$，
两实心封装表面 $\bm n\cdot\bm J_e=0$，
其余材料界面需相容的电磁跳跃条件。
以 $\bm B=\Curl\bm A$、
$\bm E=-\jj\omega\bm A-\grad\phi$ 得
\begin{equation}
\boxed{
\begin{aligned}
 \Curl(\mu^{-1}\Curl\bm A)
 &=\bm J_{\rm src}-\sigma(\jj\omega\bm A+\grad\phi),\\
 \Div[\sigma(\jj\omega\bm A+\grad\phi)]&=0
 \quad\text{（导电域）}.
\end{aligned}}
\label{eq:bfzi-direct-locked-a-phi}
\end{equation}
实际三维参考必须处理规范/零空间、拓扑、
导体--绝缘体界面和正确无穷远衰减。
人工外球的 $\bm n\cdot\bm J_e=0$
只定义另一个截断问题，不能自动等同无限海水。

\subsection{原 Green 电流弱式的条件性等价}
只在常数 $\mu=\mu_0$、合法完整源、
真正无散度零法向测试空间及收敛原核双积分
条件下，令
\[
 (\mathcal G_0\bm f)(\bm x)=
 \int_{\Omega_\infty}
 \frac{\bm f(\bm y)}{4\pi|\bm x-\bm y|}\,
 \mathrm d\bm y,\qquad
 \bm A=\mu_0\sum_{p=1}^2\bm C_pI_p+
 \mu_0\mathcal G_0\bm J_e.
\]
对实合规电流基 $\bm W_m$：
\begin{align}
 R_{mn}&=\int_{\Omega_\infty}\sigma_s^{-1}
 \bm W_m\cdot\bm W_n\,\mathrm dV,\\
 L_{mn}&=\mu_0\int_{\Omega_\infty}
 \bm W_m\cdot\mathcal G_0\bm W_n\,\mathrm dV,\\
 K_{mp}&=\mu_0\int_{\Omega_\infty}
 \bm W_m\cdot\bm C_p\,\mathrm dV .
\end{align}
由弱式得到
\begin{equation}
\boxed{
 (\bm R+\jj\omega\bm L)\bm z
 =-\jj\omega\bm K\bm I,\qquad
 \bm Z_{\rm sw}=\omega^2\bm K^T
 (\bm R+\jj\omega\bm L)^{-1}\bm K .
}
\label{eq:bfzi-direct-locked-schur}
\end{equation}
此式在已声明条件下数学上合法，
\textbf{并非当前有限 Gauss--curl 字典、
原核 1/r 求积或外球截断已经收敛的证据}。
有非均匀磁导率时，自由空间核不得直接代替
完整磁化与界面响应。

\subsection{复端口与周期平均损耗}
线性互易、无源双端口要求
\[
 \bm Z_{\rm EM}=\bm Z_{\rm EM}^T,\qquad
 \Herm{\bm Z_{\rm EM}}\succeq0,
\]
非对角互阻可取负值，不得逐项夹正。
峰值相量下必须用相同实际海水响应验算
\begin{equation}
\boxed{
 P_{\rm sea}=\frac12\int_{\Omega_\infty}
 \frac{|\bm J_e|^2}{\sigma_s}\,\mathrm dV
 =\frac12\bm I^{\mathrm H}
 \Herm{\bm Z_{\rm sw}}\bm I .
}
\end{equation}
$\bm Z_{\rm sw}$ 只表示海水增量；
铜/磁性材料/电路损耗需独立、无重数归属。
同源代数功率闭合本身不能代替
独立参考场和端口精度的证明。

\subsection{顺序验收与停止规则}
\begin{enumerate}
 \item D0：物理场景、源、MQS 适用与相量方向；
 \item D1：圆形共轴分层海水 Helmholtz 半解析与互阻；
 \item D2：独立三维 $A$--$\phi$ 求解及 TEAM-7 等
  基准、网格和远场收敛；
 \item D3：真实 BFZI 从无盒、单盒到双独立盒、
  从共轴到任意合法 SE(3)，比较两端口复阻抗和
  $\bm E,\bm B,\bm J_e,q$；
 \item D4：不同端口相位下互易、被动、
  真实铜/海水分账和端口--体积功率一致；
 \item D5：合规 BEM/FEM--BEM 或 Green--Galerkin
  与独立参考交叉验证后再择优作为高效直接法；
 \item D6：只有可信物理直接教师验收后，
  才进行空间热演化与跨几何神经代理。
\end{enumerate}
三维 FEM 允许作为\emph{独立精度参考}；
尽量避免世界体网格是后续 D5 速度目标，
不能阻止验证真值。
原 P0/P1-A/P1-B 限定范围 PASS 仍保留；
原五模态 G1/P1-C 数值工作冻结为
\textbf{未认证历史试验}，
不继续无目标调整角帽、外球及节点。

\bibliographystyle{unsrt}
\bibliography{references}
\end{document}